\textbf{РОЗДІЛ 1.}

\textbf{ТЕОРЕТИЧНІ ЗАСАДИ НАВЧАННЯ РОЗРОБКИ ІМЕРСИВНИХ ОСВІТНІХ РЕСУРСІВ}

У розділі висвітлено концептуальні основи навчання майбутніх учителів інформатики розробки імерсивних освітніх ресурсів, що постають одним із ключових напрямів педагогічної інноватики в умовах цифрової трансформації освіти. Розкрито педагогічну проблематику створення VR/AR/XR-рішень, що потребує від майбутнього вчителя інтеграції технологічних, методичних та дизайнерських компетентностей у межах цілісної професійної картини діяльності. Уточнено сутнісні характеристики імерсивних освітніх ресурсів, які представлені як інтерактивні засоби нового покоління, побудовані на моделюванні віртуальних або доповнених середовищ та спрямовані на підсилення дидактичних можливостей навчання. Здійснено систематизацію вітчизняного досвіду використання імерсивних технологій в освітньому процесі, зосереджену на аналізі освітньо-професійних програм підготовки майбутніх учителів інформатики.

\textbf{1.1. Розробка імерсивних освітніх ресурсів як педагогічна проблема підготовки майбутніх учителів інформатики}

Активне впровадження новітніх педагогічних технологій, які змінюють характер взаємодії учасників освітнього процесу, - провідна ознака часу. Освіта XXI ст. розвивається під впливом процесів цифровізації, швидкого оновлення знань і зростання ролі технологій, які забезпечують персоналізоване, адаптивне та інтерактивне навчання \hl{{[}U-1{]}}.

Посилення ролі ІКТ закономірно спрямовує увагу на імерсивні технології як засоби поглибленого занурення у навчальний зміст і розширення можливостей професійної підготовки фахівців. За таких умов вивчення імерсивних технологій набуває особливої актуальності у підготовці майбeтніх вчителів, зокрема вчителів інформатики, оскільки ІТ формують новий рівень взаємодії з навчальним середовищем, забезпечуючи занурення у зміст і розширюючи можливості педагогічної діяльності в умовах цифрової трансформації освіти. Так імерсивні освітні ресурси (далі - ІОР) стають невід'ємним складником цифрової трансформації освіти, що відбувається в умовах глобального інформаційного суспільства і формує нові вимоги до професійної діяльності педагога. Цифрова компетентність сьогодні є необхідною умовою професійної діяльності сучасного вчителя, оскільки вона є запорукою ефективного виконання вчителем педагогічних завдань та досягнення освітніх цілей (\textbf{Л.~Кібенко} {[}\hl{К-5, с. 257}{]}). Отже ``проблема породжує потребу у підготовці майбутніх учителів до формування цифрової компетентності, а особливо використання імерсивних технологій, апробацію нововведень та їх реалізацію в умовах сучасних закладів загальної середньої освіти'' \hl{{[}Ш-2, с.140{]}.}

За таких умов учитель, зокрема вчитель інформатики, має бути не лише користувачем готових цифрових рішень, а й творцем інноваційних освітніх середовищ, здатним інтегрувати технології доповненої, віртуальної та змішаної реальності у власну педагогічну практику \hl{{[}Аe-1{]}.} У підготовці майбутніх учителів інформатики розробка таких ресурсів виходить за межі суто технічного завдання, набуваючи статусу психолого-педагогічної проблеми, оскільки передбачає формування у студентів не лише цифрових навичок, а й педагогічного мислення, креативності, рефлексивності та етичної відповідальності за якість освітнього продукту \hl{{[}Ke-1{]}.}

З позицій сучасної педагогічної науки імерсивне навчання розглядається як діяльність, що створює умови для розвитку вищих психічних функцій через досвід занурення у віртуально-реальні ситуації, наближені до реальних умов майбутньої професійної діяльності. Психолого-педагогічний аспект розробки ІОР зумовлений необхідністю врахування закономірностей сприйняття, уваги, пам'яті, мотивації та емоційного залучення користувача. Від рівня усвідомлення цих процесів залежить ефективність побудови навчального сценарію, дизайну середовища та інтеграції інструментів VR/AR/MR у педагогічну практику \hl{{[}R-1{]}}.

Проблемність теми полягає у відсутності усталених методик формування готовності майбутніх учителів до створення ІОР, невизначеності критеріїв їхньої педагогічної ефективності та недостатній інтеграції психолого-педагогічних підходів до цифрового дизайну.

\hl{На цьому тлі актуалізується пошук педагогічних умов і методичних стратегій, що сприятимуть розвитку в майбутніх учителів інформатики здатності до свідомого, творчого й етично виваженого використання технологій занурення в освітньому процесі {[}\textbf{Г-1}{]}. Дослідники зазначають, що у} процесі підготовки майбутніх учителів інформатики розробка імерсивних освітніх ресурсів виступає одночасно засобом і результатом формування професійних компетентностей \hl{{[}A-1{]}.} За таких умов студент постає не лише користувачем цифрових технологій, а проєктувальником інтерактивного освітнього середовища, здатним забезпечити педагогічну доцільність, доступність і безпечність цифрового контенту. Отже, навчання розробки імерсивних освітніх ресурсів потребує цілісного психолого-педагогічного підґрунтя, що інтегрує положення когнітивної психології, педагогічного дизайну, дидактики та теорії навчання інформатики \hl{{[}D-1{]}}.

Актуальність окресленої проблеми зумовлена сукупністю соціокультурних, освітніх і технологічних трансформацій, що відбуваються у сучасному суспільстві, у зв'язку з чим вона набула статусу однієї з ключових тем науково-педагогічного дискурсу. Її обговорення у вітчизняних і зарубіжних наукових публікаціях зосереджується на переосмисленні цілей освіти, оновленні змісту підготовки фахівців і пошуку адекватних методологічних підходів до розв'язання окреслених викликів, що підтверджує теоретичну значущість і практичну затребуваність порушеної проблематики.

Аналіз джерельної бази засвідчує доцільність її систематизації за кількома укрупненими проблемно-тематичними кластерами, що відображають основні напрями науково-педагогічного осмислення імерсивних технологій в освіті (Рис.1.1.).

\includegraphics[width=6.69653in,height=3.37902in]{thesis_final/images/media/image1.png}

\emph{Рис.~1.1}. Проблемно-тематичні кластери науково-педагогічного осмислення імерсивних технологій в освіті

Перший кластер охоплює \emph{\textbf{загальнотеоретичні та дидактичні підходи до трактування імерсивних освітніх ресурсів}.} У межах цього напряму імерсивні технології розглядаються як інноваційний дидактичний інструмент, що зумовлює трансформацію змісту, форм і методів навчання, розширює можливості організації освітнього процесу та сприяє реалізації сучасних моделей навчання, орієнтованих на гнучкість і варіативність освітнього середовища.

Другий кластер репрезентує \emph{\textbf{практики впровадження імерсивних технологій у освітній процес на різних рівнях і в різних ланках освіти}.} У цьому контексті увага зосереджується на педагогічних умовах, організаційних моделях і дидактичних сценаріях використання імерсивних ресурсів з урахуванням вікових, когнітивних і освітніх особливостей здобувачів освіти.

Третій блок об'єднує \emph{\textbf{дослідження використання імерсивних технологій у системі професійної педагогічної підготовки та розвитку педагогічного професіоналізму}.} Імерсивні ресурси осмислюються як засіб формування професійних і цифрових компетентностей майбутніх учителів, розвитку їхньої методичної готовності та здатності до проєктування освітніх рішень у цифровому освітньому середовищі.

Четвертий кластер становлять \emph{\textbf{дослідження застосування імерсивних технологій в інклюзивній освіті}.} У межах цього напряму імерсивні освітні ресурси розглядаються як інструмент забезпечення доступності навчання, підтримки різноманітності освітніх потреб та створення безбар'єрного, адаптивного освітнього середовища для всіх категорій здобувачів освіти.

П'ятий кластер охоплює \emph{\textbf{галузеві, порівняльні та узагальнювальні підходи до використання імерсивних технологій в освіті}.} У цьому блоці імерсивні ресурси аналізуються як дидактичний засіб розвитку предметних і професійних компетентностей у різних освітніх галузях, а також у контексті порівняльного аналізу міжнародного й вітчизняного досвіду та узагальнення результатів наукових і прикладних досліджень.

Коротко окреслимо коло аспектів, які набули найбільшої уваги у працях вітчизняних і зарубіжних науковців та визначають провідні напрями осмислення означеної проблеми у сучасному науково-педагогічному дискурсі.

У межах \emph{\textbf{першого кластера наукових джерел}} актуалізується проблематика класифікації імерсійних технологій, визначення їх дидактичного потенціалу та окреслення міждисциплінарних можливостей застосування в освітньому процесі. Дослідники зосереджуються на осмисленні VR/AR-технологій як інноваційного дидактичного інструментарію, здатного трансформувати зміст навчання, розширити спектр освітніх сценаріїв і підсилити пізнавальну активність здобувачів освіти.

Значний внесок у розроблення класифікаційних підходів до імерсивних технологій здійснено у працях Ю.Носенко (2024) \hl{{[}Н-2{]},} де запропоновано структурну типологію імерсивних сервісів для освітнього процесу. Авторка систематизує занурювальні середовища за низкою критеріїв, серед яких провідними виступають педагогічна мета використання, рівень інтерактивності, спосіб візуалізації та технічна складність. У роботі наголошується, що результативність упровадження імерсивних технологій визначається не лише технічними характеристиками, а й узгодженістю цифрових інструментів з дидактичними завданнями та готовністю користувачів до взаємодії з віртуальними моделями. Проблема \emph{персоналізації навчання в умовах цифрового середовища} розкривається у дослідженні О.Гурської та О.Самборської (2025) \hl{{[}Г-10{]},} де VR/AR-технології розглядаються як засіб побудови індивідуальних освітніх маршрутів. Авторки акцентують увагу на можливостях автоматизованого аналізу навчального прогресу здобувачів освіти, що створює підґрунтя для адаптації навчального контенту відповідно до індивідуальних освітніх потреб і темпу засвоєння матеріалу.

Окремий напрям у межах цього кластера становлять праці, присвячені \emph{дидактичному потенціалу доповненої реальності у шкільній освіті.} Зокрема, О. Пінчук і Л.Лупаренко (2022) \hl{{[}П-3{]}} доводять, що AR-контент забезпечує достовірне моделювання природничих процесів, поєднуючи візуалізацію з дослідницькою діяльністю учнів. Автори підкреслюють можливість використання доповненої реальності як засобу активізації навчального пізнання та формування дослідницьких умінь, наводячи приклади цифрових ресурсів, адаптованих до потреб шкільного навчання.

Розвиток ідей синтетичного та гібридного освітнього середовища послідовно представлено у працях Ю. Богачкова (2023) \hl{{[}Б-5{]}.} Автор вводить поняття «імерсивний синтетичний простір», у межах якого поєднуються очні й дистанційні форми навчання. У дослідженнях підкреслюється, що такі середовища сприяють глибшому зануренню в навчальну діяльність, забезпечують гнучкість освітніх сценаріїв і розширюють можливості індивідуального та групового навчання.

Психолого-педагогічний вимір застосування імерсивних технологій розкрито у дослідженні С.Литвинової (2025) \hl{{[}Л-4{]}}, де обґрунтовується потенціал VR/AR-середовищ для розвитку емоційного інтелекту учнів. Згідно з отриманими результатами, імерсивні навчальні ситуації створюють умови для усвідомлення, регуляції та конструктивного вираження емоцій, що відкриває перспективи формування психологічно орієнтованої цифрової педагогіки.

У контексті проєктування електронних освітніх ресурсів О.Чорноус (2025) \hl{{[}Ч-2{]}} звертає увагу на здатність імерсивних технологій посилювати ефект присутності, покращувати запам'ятовування навчального матеріалу та формувати відчуття «дії від першої особи». Дослідниця наголошує на доцільності інтеграції VR-практик у систему цифрової дидактики як складника сучасного освітнього дизайну. Подібну проблематику порушує й О.Орлов (2024) \hl{{[}О-2{]},} який розглядає імерсивні технології як елемент інноваційних освітніх маршрутів, що поєднують інтерактивність, візуальну насиченість і варіативність навчання. Автор акцентує на необхідності методичної підтримки педагогів для ефективного використання VR/AR-інструментів у практиці навчання.

Праці першого кластера формують цілісне уявлення про імерсивні технології як багатовимірний дидактичний феномен, у межах якого поєднуються класифікаційні підходи, моделювання навчальних середовищ, персоналізація освітніх траєкторій і психолого-педагогічні аспекти навчання, що створює методологічне підґрунтя для подальшого аналізу їх практичної реалізації в освітньому процесі різних рівнів і ланок освіти. Саме на цих засадах у працях другого кластера зосереджується увага на конкретних практиках упровадження імерсивних технологій у систему освіти, педагогічних умовах їх використання та дидактичних сценаріях застосування в реальних освітніх контекстах (другий кластер).

У межах \emph{другого кластера} наукових джерел у центрі уваги перебуває проблематика практичного впровадження імерсивних технологій в освітній процес на різних рівнях і в різних ланках освіти, а також визначення педагогічних, методичних і технічних умов їх ефективного використання. Дослідження цього напряму репрезентують перехід від теоретичного осмислення імерсивних ресурсів до аналізу реальних освітніх практик, що охоплюють початкову, базову, профільну середню та вищу освіту, а також різні освітні галузі. При цьому звертається увага не лише на позитивні результати використання ІОР, а й на проблемні аспекти.

Застосування імерсивних технологій у профільній середній освіті розглядається у працях О. Омельчук зі співавторами (2024) \hl{{[}О-1{]},} де акцент зроблено на їхньому потенціалі для розвитку просторового мислення, творчості та технологічної грамотності учнів. Автори підкреслюють, що використання віртуальних лабораторій, симуляцій і тривимірних моделей у технологічній освітній галузі посилює практичну спрямованість навчання та підвищує мотивацію до пізнавальної діяльності.

Проблеми інтеграції імерсивних ресурсів у початкову освіту висвітлено у дослідженні М. Швардак (2023) \hl{{[}Ш-1{]}}, де імерсивні технології розглядаються як складник ігрової педагогіки. Авторка доводить, що такі засоби сприяють підтриманню навчального інтересу молодших школярів, розвитку емоційного інтелекту та формуванню базових цифрових навичок через поєднання пізнавальної діяльності з емоційним залученням. Використання VR/AR-технологій як інструменту підтримки змішаного навчання у закладах загальної середньої освіти обґрунтовується у праці Ю. Носенко \hl{{[}Н-6{]}.} Авторка зазначає, що імерсивні рішення відповідають концептуальним засадам «Нової української школи», оскільки сприяють індивідуалізації навчання, підвищенню мотивації учнів і візуалізації складних понять, зокрема в межах природничо-наукових дисциплін.

Технічний і методичний вимір упровадження імерсивних технологій розкривається у публікації Т. Цимбалюк і Д. Федасюка (2024) \hl{{[}Ц-1{]},} де здійснено класифікацію освітніх імерсивних середовищ, визначено переваги VR-платформ порівняно з традиційними мультимедійними засобами та окреслено вимоги до апаратного забезпечення. Автори наголошують на необхідності добору рівня занурення з урахуванням вікових і когнітивних особливостей здобувачів освіти.

Вагомий внесок у формування міждисциплінарного наукового дискурсу щодо освітнього потенціалу імерсивних технологій здійснили Ю. Носенко та А. Сухіх (2025) \hl{{[}Н-1{]},} які узагальнили результати Міжнародної конференції ITE-2025. У їхніх матеріалах VR/AR-технології постають як стратегічний ресурс розбудови інноваційних педагогічних моделей і цифрових освітніх екосистем.

Практичні аспекти використання імерсивних технологій у STEM- та STEAM-освіті на рівні початкової школи розглядаються у дослідженні А. Дрокіної (2025) \hl{{[}Д-3{]}.} Авторка доводить, що VR/AR-засоби сприяють формуванню дослідницьких навичок, просторового мислення та інтересу до природничо-технічних дисциплін, водночас підкреслюючи важливість створення доступних освітніх платформ для дітей різного віку.

Окрему групу праць другого кластера становлять дослідження, присвячені використанню імерсивних технологій в умовах кризових і складних соціальних обставин. Зокрема у роботах І. Удовиченко та В. Удовиченко (2024) \hl{{[}У-2{]}} обґрунтовано теоретико-методологічні засади застосування VR/AR-засобів у закладах загальної середньої освіти в умовах воєнного стану. Авторки розглядають імерсивні технології як інструмент забезпечення безперервності навчання, організації дистанційних і змішаних форматів, підтримки психологічної стабільності учнів і реалізації принципів рівного доступу до якісної освіти. У працях систематизовано види імерсивних технологій, окреслено критерії добору цифрових інструментів і подано приклади використання віртуальних лабораторій та симуляцій у навчальному процесі.

Проблеми організації шкільного експерименту з використанням VR/AR-технологій проаналізовано у дослідженні О. О. Соколюк (2024) {[}С-5{]}, де доведено ефективність віртуальних лабораторій у навчанні фізики й астрономії з погляду підвищення якості засвоєння складних понять і дотримання принципів безпечного навчання.

Методологічні засади створення імерсивних освітніх ресурсів розглянуто у колективній праці С. Семерікова зі співавторами (2022) \hl{{[}S-1{]},} у якій запропоновано дизайн-модель розроблення VR/AR-засобів навчання з урахуванням дидактичних принципів, когнітивного навантаження та оцінювання ефективності занурювального досвіду. Автори наголошують на важливості педагогічного дизайну й комунікативної взаємодії між учасниками освітнього процесу.

У контексті вищої освіти К. Гнедіна (2025) \hl{{[}Г-11{]}} аналізує імерсивні технології як чинник забезпечення якості навчання, підкреслюючи їхній вплив на залученість студентів, розвиток критичного мислення та академічної мотивації, а також звертаючи увагу на проблеми інтеграції VR/AR у системи моніторингу результатів навчання. Подібні ідеї простежуються у працях Ю. Матвієнка (2023) \hl{{[}М-3{]}}, де імерсивні технології розглядаються як засіб модернізації освітнього середовища, підвищення інтерактивності навчання та забезпечення доступності освітніх матеріалів через відкриті цифрові платформи.

Стратегічний вимір упровадження імерсивних технологій у діяльність освітніх закладів окреслено у дослідженні О. Кінджибали зі співавторами (2025) \hl{{[}К-10{]},} де VR/AR-засоби позиціонуються як інструмент креативного навчання, спрямованого на формування практикоорієнтованих компетентностей, критичного мислення та рефлексії. Завершує спектр праць другого кластера дослідження М. Медведєвої (2024) \hl{{[}М-4{]},} у якому імерсивні технології в STEAM-освіті визначено як каталізатор інтеграції природничих, технічних, мистецьких і гуманітарних знань, що сприяє розвитку інженерного мислення, креативності та міждисциплінарного світогляду учнів.

Іноземні вчені (Паронг, Дж., Майєр, Р. Е. \hl{{[}Ре-2{]})} аналізують \textbf{вплив ІОР на навчальні результати}, а також які \textbf{афективні та когнітивні процеси} опосередковують цей вплив.

Праці другого кластера засвідчують поступове зміщення наукового інтересу від концептуального осмислення імерсивних технологій до аналізу реальних освітніх практик їх упровадження на різних рівнях освіти. Дослідження цього напряму демонструють різноманіття педагогічних сценаріїв використання VR/AR-ресурсів - від початкової й профільної школи до вищої освіти, у межах традиційного, змішаного та дистанційного навчання, а також в умовах соціальних і безпекових викликів. Узагальнені результати дозволяють констатувати, що ефективність імерсивних технологій у навчанні зумовлюється поєднанням дидактичної доцільності, технічної готовності освітнього середовища та методичної підготовленості педагогів.

На цьому підґрунті у \emph{працях третього кластера} фокус наукового аналізу зміщується на \textbf{питання професійної педагогічної підготовки та розвитку педагогічного професіоналізму,} у межах яких імерсивні технології розглядаються вже не лише як засіб навчання учнів, а як інструмент формування професійних, цифрових і методичних компетентностей майбутніх учителів. Дослідження цього напряму фокусуються на визначенні педагогічних, організаційних і методичних умов формування готовності майбутніх учителів до проєктування та використання VR/AR/MR-технологій у професійній діяльності, а також на переосмисленні змісту й структури педагогічної освіти в умовах цифрової трансформації.

Питання формування технологічної готовності майбутніх педагогів ґрунтовно розглянуто у праці Л. Шевчук та А. Григоренка \hl{{[}Ш-4{]}.} Автори трактують технологічну готовність як інтегральну характеристику особистості педагога, що поєднує когнітивний, операційний, мотиваційний і рефлексивний компоненти та забезпечує здатність ефективно інтегрувати імерсивні технології в освітній процес. У роботі систематизовано підходи до класифікації VR/AR/MR-технологій, окреслено їхній дидактичний потенціал для створення інтерактивного освітнього середовища, розвитку креативного мислення та підвищення навчальної мотивації. Особливу увагу приділено педагогічним умовам підготовки майбутніх учителів, зокрема інтеграції спеціалізованих дисциплін у зміст педагогічної освіти, створенню практикоорієнтованого освітнього середовища та розвитку професійної рефлексії.

У контексті модернізації програм педагогічної освіти та розвитку дистанційних форматів підготовки проблему використання імерсивних технологій розглядає \textbf{Н.Яремчук \hl{{[}Я-2{]}.}} Авторка концептуалізує VR/AR/MR-засоби як новітню форму навчальної наочності, що базується на ефекті занурення та суттєво змінює сприймання навчального образу, розширюючи можливості візуалізації освітнього контенту. У дослідженні проаналізовано програмно-технічний потенціал імерсивних технологій, їхні дидактичні можливості, переваги й обмеження у професійній підготовці вчителів початкової школи, а також обґрунтовано організаційно-педагогічні, методичні та психологічні чинники ефективного використання таких інструментів у цифровому освітньому середовищі.

Окрему групу праць третього кластера становлять дослідження, у яких імерсивні технології розглядаються крізь призму галузевої підготовки педагогів. Так, О. Кириєнко (2023) \hl{{[}К-6{]}} аналізує процес формування професійної компетентності майбутніх учителів біології засобами VR/AR-технологій. Автор наголошує на можливостях моделювання складних біологічних процесів, що важко відтворюються в реальних умовах, та підкреслює значення імерсивних ресурсів для розвитку дослідницьких умінь студентів і підвищення якості природничої підготовки.

Роль імерсивних технологій у формуванні підприємницької компетентності вчителя розкривається у праці О. Кравчиної (2023) \hl{{[}К-7{]}}. Авторка доводить, що цифрові симуляції, бізнес-квести та віртуальні навчальні середовища створюють умови для розвитку критичного мислення, командної взаємодії та креативності, які є базовими складниками підприємницької освіти в умовах цифрової економіки.

Синергетичний потенціал доповненої, віртуальної та змішаної реальності в системі підготовки педагогів обґрунтовано у дослідженні Я. Ангелова (2025) \hl{{[}А-1{]}.} Автор розглядає поєднання VR, AR і MR як основу нової педагогічної парадигми, у межах якої навчання набуває адаптивного, інтерактивного й занурювального характеру. Особливий акцент зроблено на необхідності формування цифрової культури викладачів педагогічних університетів і їхньої готовності до проєктування імерсивних освітніх курсів. У ширшому контексті розвитку екосистеми освітніх технологій імерсивні рішення розглядаються у праці О. Юденкової, З. Височан і О. Сопіної (2023) \hl{{[}Ю-1{]}.} Авторки трактують VR/AR-технології як складову EdTech-екосистеми, що забезпечує цифрову стійкість і неперервність освіти, підкреслюючи їхню значущість для підготовки майбутніх педагогів в умовах трансформації освітніх систем.

Отже, праці третього кластера засвідчують, що імерсивні технології в сучасному науково-педагогічному дискурсі дедалі частіше осмислюються як інструмент трансформації системи професійної педагогічної підготовки та розвитку педагогічного професіоналізму. Узагальнення результатів досліджень дозволяє констатувати зміщення акцентів від поодинокого використання VR/AR/MR-засобів до цілісного формування технологічної, методичної та цифрової готовності майбутніх учителів, а також до переосмислення змісту й організації педагогічної освіти в умовах цифрової трансформації та розвитку EdTech-екосистем. На цьому тлі у працях четвертого кластера фокус наукового аналізу закономірно спрямовується на проблематику використання імерсивних технологій в інклюзивній освіті, де питання доступності, адаптивності та безбар'єрності освітнього середовища набувають визначального значення.

У межах четвертого кластера наукових джерел у центрі уваги перебуває проблематика використання імерсивних освітніх ресурсів в умовах інклюзивної парадигми освіти. Дослідження цього напряму зосереджені на осмисленні VR/AR-технологій як інструментів забезпечення доступності навчального контенту, компенсації сенсорних, когнітивних і комунікативних обмежень, а також створення безпечного, адаптивного й варіативного освітнього середовища для здобувачів освіти з різними освітніми потребами. У цьому контексті імерсивні технології розглядаються не лише як дидактичне нововведення, а як механізм реалізації принципів безбар'єрності, універсального дизайну навчання та рівного доступу до якісної освіти.

Можливості застосування імерсивних технологій у системі інклюзивної освіти розкриває Ю. Чигріна (2023) \hl{{[}Ч-3{]},} яка на основі аналізу зарубіжних наукових досліджень обґрунтовує ефективність VR-інструментів у роботі з учнями з особливими освітніми потребами. Авторка підкреслює, що імерсивні середовища дозволяють мінімізувати комунікативні бар'єри, знизити рівень тривожності та створити безпечний навчальний простір, адаптований до індивідуальних можливостей і потреб здобувачів освіти.

Практичний вимір використання імерсивних технологій в інклюзивному навчанні представлено у дослідженні Ю. Матвієнка, С. Нестулі та С. Шари (2022) \hl{{[}М-5{]},} де описано досвід застосування занурювальних ресурсів у роботі з дітьми з особливими освітніми потребами. Автори зазначають, що VR/AR-засоби полегшують сприйняття навчальної інформації, сприяють соціальній інтеграції та створюють умови для використання імерсивних технологій як елементів розвивальної й корекційної підтримки. Проблема індивідуалізації та соціалізації здобувачів освіти з ООП у межах змішаного навчання розглядається у праці А. Сухіх і І. Полященко (2024) \hl{{[}С-2{]}.} Авторки доводять, що використання VR-платформ у поєднанні з традиційними та дистанційними формами навчання розширює можливості включення дітей з особливими освітніми потребами у спільну навчальну діяльність, сприяє розвитку комунікативних навичок і забезпечує адаптацію освітнього процесу до індивідуальних освітніх траєкторій.

У новітніх дослідженнях українських науковців інклюзивний потенціал імерсивних технологій дедалі частіше розглядається у взаємозв'язку з іншими провідними освітніми напрямами, зокрема STEM-освітою, гуманітарною підготовкою та емоційно-ціннісним розвитком. Такий міждисциплінарний підхід засвідчує розширення наукового бачення імерсивних технологій як комплексного ресурсу, здатного одночасно підтримувати пізнавальний, соціальний і особистісний розвиток здобувачів освіти в умовах інклюзивного освітнього середовища.

Узагальнення наукових позицій дозволяє констатувати, що VR/AR-ресурси в умовах інклюзивної освіти виконують не лише компенсаторну функцію, а й виступають засобом розширення освітніх можливостей, соціальної взаємодії та особистісного розвитку здобувачів освіти. Вони забезпечують реалізацію принципів безбар'єрності, універсального дизайну навчання й рівного доступу до якісної освіти, що істотно посилює гуманістичний потенціал сучасного освітнього середовища. На цьому підґрунті у працях п'ятого кластера науковий аналіз зосереджується на галузевих, порівняльних і узагальнювальних підходах до використання імерсивних технологій в освіті, що дає змогу співвіднести національний і міжнародний досвід, а також виявити спільні тенденції та специфічні особливості розвитку цього напряму.

\emph{У межах п'ятого кластера наукових} джерел зосереджено галузеві, порівняльні та узагальнювальні підходи до використання імерсивних технологій в освіті, що дозволяють простежити еволюцію цього напряму в міжнародному та вітчизняному контекстах, а також виявити міждисциплінарні тенденції його розвитку. Дослідження цього кластера репрезентують імерсивні технології як стратегічний інструмент цифрової трансформації освіти, здатний інтегрувати педагогічні, технологічні й гуманітарні виміри навчання.

Значне місце в межах кластера посідають праці, присвячені узагальненню зарубіжного досвіду використання імерсивних технологій. Так, І. Малицька (2022) \hl{{[}М-2{]}} аналізує практики країн Європейського Союзу, зокрема Великої Британії, Німеччини та Скандинавських держав, де VR- і AR-технології системно інтегруються в підготовку педагогів, розвиток цифрової грамотності та STEM-освіту. Авторка підкреслює, що результативність упровадження імерсивних рішень у цих країнах забезпечується узгодженістю державної освітньої політики, стандартів професійного розвитку вчителів і цифрової інфраструктури.

Порівняльний аналіз зарубіжних і українських тенденцій розвитку імерсивного навчання здійснено у працях С.Литвинової (2023) \hl{{[}Л-8{]}}, де акцентовано увагу на необхідності комплексного підходу до впровадження VR/AR-технологій. Дослідниця визначає імерсивні технології як стратегічний напрям цифровізації освіти, що потребує поєднання інфраструктурного забезпечення, системної підготовки педагогічних кадрів і нормативно-правової підтримки. Подібні ідеї розвиває А. Григоренко (2024) \hl{{[}Г-3{]}}, який узагальнює вітчизняний і зарубіжний досвід підготовки педагогів із використанням імерсивних технологій, підкреслюючи, що в країнах ЄС такі інструменти вже інтегровані до стандартів педагогічної майстерності, тоді як в Україні вони перебувають на етапі становлення, зберігаючи значний потенціал професійного зростання вчителя.

Окрему групу праць п'ятого кластера становлять узагальнювальні та аналітичні видання, спрямовані на систематизацію наукового доробку та освітніх практик. У бібліографічному покажчику С. Литвинової, Ю. Носенко та В. Новика (2024) \hl{{[}Л-2{]}} здійснено огляд публікацій науковців НАПН України, що засвідчує стійку тенденцію до інтеграції імерсивних технологій у систему змішаного й дистанційного навчання. Комплексні аналітичні матеріали О. Бурова, С. Литвинової, Ю. Носенко та А. Сухіх (2024) \hl{{[}А-3{]}} обґрунтовують можливості створення інтерактивних освітніх просторів у закладах загальної середньої освіти, наголошуючи на важливості міждисциплінарної співпраці педагогів, методистів і розробників цифрових платформ. У колективній монографії С. Литвинової, Н. Сороко та С. Баценко (2023) \hl{{[}П-11{]}} запропоновано модель проєктування освітнього середовища з використанням VR/AR-засобів, ефективність якої пов'язується з дотриманням принципів активності, безпечності та динамічної педагогічної взаємодії.

Важливим складником п'ятого кластера є дослідження галузевого та міждисциплінарного застосування імерсивних технологій. Так, Т. Гранчак (2023) \hl{{[}Г-6{]}} аналізує використання імерсивних продуктів у сфері документування та популяризації історико-культурної спадщини, підкреслюючи потенціал віртуальних музеїв і 3D-експозицій для бібліотечної та інформаційної практики. О. Бойко (2024) \hl{{[}Б-2{]}} доводить ефективність імерсивних середовищ у навчанні англомовної письмової комунікації, акцентуючи на розвитку навичок саморегуляції та контекстного розуміння. Застосування VR-технологій у поліграфічній освіті розглядають А. Гордєєв і П. Чернікова (2025) \hl{{[}Г-9{]},} які наголошують на формуванні дизайнерського мислення й нових професійних компетентностей.

Розширення гуманітарного виміру імерсивних технологій простежується у працях С. Доценка і В. Чжен (2023) \hl{{[}Д-1{]},} де описано феномен «імерсивного шоу» як синтез мистецтва, театру та цифрових технологій, а також у дослідженнях мовної освіти. Зокрема, О. Кулик (2024) \hl{{[}К-9{]}} аналізує використання імерсивних платформ у навчанні української мови, підкреслюючи їхню роль у подоланні комунікативних бар'єрів, тоді як Н. Сороко та О. Гаєвська (2021) \hl{{[}С-7{]}} обґрунтовують ефективність VR-модулів у викладанні східних мов. Питання іншомовної підготовки також розкрито у праці К.Бровко й О.Лободи (2024) \hl{{[}Б-6{]},} де імерсивне навчання розглядається як цілісний дидактичний процес, спрямований на розвиток критичного мислення та саморегуляції.

У контексті технологічного розвитку професійної освіти Н. Чубінська та Д. Попадюк (2025) \hl{{[}Ч-1{]}} простежують взаємозв'язок між імерсивними технологіями й штучним інтелектом, акцентуючи на їхньому значенні для розвитку адаптивного навчання та підвищення інклюзивності освітніх систем. Подібні техніко-організаційні аспекти впровадження VR/AR у змішаному навчанні окреслює С.Баценко (2025) \hl{{[}Б-3{]},} наголошуючи на ролі інфраструктури, фінансового забезпечення та кадрової готовності педагогів. Розвиток гуманітарних дисциплін у цифровому просторі доповнюють ідеї Л.Біловус і О. Васильків (2024) \hl{{[}Б-4{]}}, які пропонують концепцію «імерсивного мультифункціонального хабу» як середовища міжкультурної комунікації та інтерпретації культурного контенту.

Вагомий сегмент п'ятого кластера становлять дослідження галузевого застосування імерсивних технологій у мовній, гуманітарній та технічній освіті, що засвідчує розширення міждисциплінарного простору їх використання. Зокрема, О. Кулик (2024) \hl{{[}К-9{]}} досліджує роль імерсивних технологій у навчанні української мови, акцентуючи на використанні цифрових платформ для імітації мовленнєвих ситуацій, діалогів і проєктної діяльності, що посилює комунікативну складову мовної освіти. Авторка зазначає, що імерсивні ресурси сприяють подоланню психологічних бар'єрів у спілкуванні та формуванню мовної впевненості здобувачів освіти.

У праці І. Сабодашко й Т. Горохівської (2024) \hl{{[}С-6{]}} проаналізовано взаємозв'язок між імерсивними технологіями та штучним інтелектом у гуманітарній освіті. Авторки наголошують, що VR/AR-засоби не замінюють традиційне навчання, але виступають його важливим доповненням, особливо у процесі формування медіа- й цифрової грамотності, а також критичного осмислення навчального контенту. Галузевий вимір імерсивного навчання у технічній освіті розкрито у дослідженні Л. Нікітіної та Н. Дженіюк (2023) \hl{{[}Н-3{]},} де імерсивні технології розглядаються як ефективний інструмент підготовки фахівців у сфері телекомунікацій. Авторки доводять, що інтерактивні 3D-моделі та графічні симуляції сприяють глибшому засвоєнню теоретичних положень дисциплін, розвитку просторового мислення та формуванню навичок технічного аналізу. Завершує наш огляд досліджень п'ятого кластера праця С. С. Ізбаша (2020) \hl{{[}І-3{]},} у якій імерсивні технології розглядаються як складова відкритих освітніх ресурсів, здатна посилювати не лише доступність навчання, а й активну участь здобувачів освіти у пізнавальному процесі, що логічно поєднує ідеї відкритої освіти з сучасними цифровими інноваціями.

Загалом праці п'ятого кластера репрезентують імерсивні технології як міждисциплінарний і стратегічний ресурс розвитку сучасної освіти, що виходить за межі окремих рівнів і галузей навчання. Узагальнення галузевого, порівняльного та аналітичного досвіду засвідчує, що ефективність упровадження VR/AR-рішень зумовлюється не лише технічними інноваціями, а й системністю освітньої політики, підготовкою педагогічних кадрів, міждисциплінарною взаємодією та орієнтацією на гуманістичні цінності освіти. У цьому контексті імерсивні технології постають як інструмент інтеграції педагогічних, культурних, технологічних і соціальних вимірів освітнього процесу.

З метою систематизації наукових підходів до вивчення імерсивних освітніх ресурсів та виявлення логіки їх проблемно-тематичного осмислення в сучасному науково-педагогічному дискурсі доцільно представити деталізовану модель кластерного групування досліджень, що конкретизує Рис.1.2.

\includegraphics[width=6.69653in,height=6.43113in]{thesis_final/images/media/image2.png}

\emph{Рис.1.2.} Змістове наповнення проблемно-тематичних кластерів науково-педагогічного осмислення імерсивних технологій в освіті.

Подана модель відображає взаємопов'язаність проблемно-тематичних кластерів імерсивних технологій в освіті та демонструє інклюзивний освітній вимір як наскрізну точку перетину дидактичних, практичних, професійно-педагогічних і галузево-порівняльних підходів.

Оглянуті праці засвідчують узгодженість позицій дослідників щодо зростання ролі імерсивних технологій у сучасному освітньому просторі, водночас у площині педагогічної освіти акценти зміщуються від фрагментарного використання цифрових інструментів до їх системного впровадження у професійну підготовку вчителя. Дослідження фіксують посилення наукового інтересу до проблеми практичної реалізації імерсивних технологій у різних сегментах освіти~-- від загальної середньої школи до системи професійної педагогічної підготовки. Новітні публікації також засвідчують розширення міждисциплінарного простору їх застосування, зокрема в мистецькій, лінгвістичній, економічній та гуманітарній освіті, що виходить за межі суто шкільного або педагогічного контексту. У науковому дискурсі імерсивні технології дедалі частіше осмислюються не лише як інноваційні дидактичні інструменти, а як системоутворювальний чинник підготовки педагога цифрової епохи, зорієнтованого на інтеграцію когнітивних, емоційних і практичних компонентів навчальної діяльності.

\textbf{1.2. Сутнісна характеристика та структура процесу навчання розробки імерсивних освітніх ресурсів}

Трансформації освітнього середовища, зумовлені цифровізацією, поширенням технологій занурення та зміною ролі вчителя в умовах Нової української школи, актуалізують переосмислення процесу професійної підготовки майбутніх учителів інформатики. Аналітичне вивчення наукової літератури (підрозділ 1.1.) свідчить про те, що процес навчання розробки імерсивних освітніх ресурсів доцільно розглядати не лише як сукупність навчальних дій і технологічних операцій, а як цілісну систему, що інтегрує мотиваційні, когнітивні, діяльнісні та рефлексивні компоненти професійної підготовки майбутнього вчителя. Сутність процесу визначається взаємодією педагогічних, психологічних і технологічних чинників, а також орієнтацією майбутнього вчителя на досягнення освітньо значущих результатів, пов'язаних із проєктуванням безпечних, доступних і педагогічно доцільних імерсивних ресурсів.

З огляду на викладене виникає потреба в окресленні структурної організації означеного процесу, визначенні його ключових складників, логіки їхньої взаємодії та функціонального призначення в системі професійної підготовки майбутніх учителів інформатики. Передусім доцільно окреслити термінологічний апарат дослідження, оскільки процес навчання розробки імерсивних освітніх ресурсів пов'язаний із використанням значної кількості міждисциплінарних понять, що функціонують у педагогічному, психологічному та технологічному дискурсах і не завжди мають усталене трактування в межах професійної педагогічної освіти.

У межах дослідження суттєвого значення набувають поняття, що безпосередньо відображають специфіку імерсивних технологій як об'єкта педагогічного проєктування, роль цифрового дизайну в освітньому процесі та професійну підготовку майбутніх учителів інформатики. До таких базових понять належать: ``імерсивні технології'' (далі - ІТ) як сукупність цифрових засобів і середовищ занурення; ``імерсивні освітні ресурси'' (далі - ІОР) як педагогічно спроєктовані цифрові продукти; процес навчання розробки імерсивних освітніх ресурсів (далі - ПНР\_ІОР); цифровий дизайн (далі - ЦД) в освіті та його диференційовані варіанти (digital instructional design, digital learning design, digital media design); професійна готовність майбутнього вчителя інформатики до розробки імерсивних освітніх ресурсів; педагогічна ефективність імерсивних освітніх ресурсів. Окреслені поняття потребують детальнішого розгляду з метою конкретизації їх змісту та ролі у структурі досліджуваного процесу.

Передусім зазначимо, що базове поняття ``імерсивність'' використовується у різних аспектах та переважно визначається як занурення, тобто стан глибокої когнітивної, емоційної та сенсорної залученості суб'єкта в певну діяльність або штучно створене середовище, що забезпечує ефект присутності та інтенсифікує взаємодію з навчальним або віртуальним простором \hl{{[}Ш-3, с. 14{]}.}

\textbf{Імерсивні технології} у сучасному науково-педагогічному дискурсі розглядаються як сукупність цифрових засобів і середовищ, що забезпечують повне або часткове занурення суб'єкта навчання у віртуально сконструйований простір. У науковій літературі зазначені технології тлумачаться як інструменти створення ефекту присутності та взаємодії людини з цифровим середовищем, що передбачає її залучення до штучно змодельованого або змішаного простору \hl{{[}С-13{]}.} Такий підхід поділяють і зарубіжні науковці, зокрема \textbf{Й. Радіанті (J. Radianti)} та співавтори, \textbf{Дж. Паронг (J. Parong)} і \textbf{Р. Е. Маєр (R. E. Mayer)}, \textbf{Г. Макранскі (G. Makransky)}, у працях яких імерсивність (\emph{immersion}) трактується передусім як властивість технологічного середовища навчання, зумовлена рівнем технічного занурення та формуванням відчуття присутності \emph{(presence)}, а не як самостійно концептуалізований освітній ресурс \hl{{[}Ре-2; Re-2{]}. Водночас у зарубіжному науковому дискурсі представлена й інша дослідницька позиція, відповідно до якої імерсивність не зводиться виключно до технічних характеристик середовища або використання спеціалізованих пристроїв занурення. Так, у працях \textbf{Дж. Г. Робертсона (G. G. Robertson)}, \textbf{М. Червінські (M. Czerwinski)} та \textbf{М. ван Данціха (M. van Dantzich)} імерсивність розглядається як психологічний ефект присутності, що може формуватися і в настільних (\emph{desktop}) віртуальних середовищах за умови відповідного дизайну взаємодії, інтерфейсу та організації уваги користувача. У такому підході імерсивність постає не як похідна від рівня сенсорного занурення, а як результат цілісної взаємодії користувача з цифровим середовищем, що розширює межі розуміння імерсивних технологій поза рамками виключно повністю імерсивної віртуальної реальності {[}Re-3{]}.}

\hl{Також звертаємо увагу на те, що в зарубіжному науковому дискурсі імерсивність та присутність не ототожнюються автоматично з використанням повністю імерсивних пристроїв, а розглядаються крізь призму когнітивного та психологічного досвіду користувача. У статті \textbf{Е. Крокоса (Е.Krokos), К. Плезан (С.Plaisant) та А. Варшні} \textbf{(А.Varshney)} імерсивність (\emph{immersion}) пов'язується насамперед із можливостями віртуального середовища залучати просторову пам'ять і сенсомоторні механізми людини. Автори зазначають, що віртуальні середовища «можуть поєднувати візуально імерсивні просторові репрезентації даних із нашими вестибулярними та пропріоцептивними відчуттями» {[}Ke-2, с.1{]}, підкреслюючи роль тілесного досвіду у формуванні когнітивних процесів.}

У розширеному трактуванні імерсивні технології охоплюють різні форми віртуальної, доповненої та змішаної реальності, які забезпечують як повне занурення у віртуальний світ, так і інтеграцію цифрових об'єктів у реальне середовище \hl{{[}D-2; Ш-2, с. 141{]}. Зважаючи на існування реальної (об'єктивної) реальності \emph{(RR -- real reality)} ``в якій знаходиться користувач і яку сприймає органами своїх чуттів. Це світ, що існує незалежно від нього і його свідомості'' {[}Ш-1, с.230{]}, на думку вчених,} головна ідея імерсивних технологій полягає у підвищенні відчуття присутності, взаємодії та участі користувача у віртуальному або поєднаному (реальному та віртуальному) середовищі \hl{{[}Ш-1, с.230; С-14{]}.}

\emph{\textbf{Доповнена реальність} (AR} -- \emph{augmented reality)} -- технологія, яка доповнює реальний світ віртуальними елементами, накладаючи їх на зображення реального світу, яке бачить користувач. AR дозволяє користувачам взаємодіяти з цими віртуальними об\textquotesingle єктами в реальному часі. Серед основних характеристик: доповнення реального світу (AR не замінює реальний світ, а лише додає до нього віртуальні елементи); взаємодія в реальному часі (користувачі можуть взаємодіяти з віртуальними об\textquotesingle єктами в реальному часі, змінюючи їх положення, розмір або інші параметри); використання різних пристроїв (AR може бути реалізована на різних пристроях, таких як смартфони, планшети, окуляри AR та інші). Застосунки AR працюють із тривимірними моделями, текстами, зображеннями, анімацією та відео \hl{{[}Т-2{]}}

\emph{\textbf{Віртуальна реальність} (VR} -- \emph{virtual reality)} -- визначається як технологія, яка створює повністю змодельовані середовища, які занурюють користувача в новий світ \hl{{[}К-17{]}}. VR використовує спеціальні пристрої, такі як VR-шоломи, для створення ілюзії присутності в віртуальному світі. До основних характеристик VR віднесено: повне занурення (VR повністю замінює реальний світ віртуальним, створюючи ілюзію присутності в іншому місці); інтерактивність (користувачі можуть взаємодіяти з віртуальним світом, використовуючи контролери або інші пристрої введення); використання VR-шоломів (VR зазвичай вимагає використання VR-шоломів, які закривають очі користувача та відображають віртуальне зображення).

\emph{Змішана реальність (MR - mixed reality)} -- технологія, яка комбінує елементи як реального, так і віртуального світу. MR дозволяє віртуальним об\textquotesingle єктам взаємодіяти з реальним світом, створюючи гібридне середовище. Серед основних характеристик: комбінація реального та віртуального (MR поєднує елементи реального та віртуального світу, створюючи гібридне середовище); взаємодія віртуальних об\textquotesingle єктів з реальним світом (віртуальні об\textquotesingle єкти в MR можуть взаємодіяти з реальним світом, наприклад, відбиватися від реальних поверхонь або реагувати на реальні об\textquotesingle єкти); використання спеціальних пристроїв (MR зазвичай вимагає використання спеціальних пристроїв, таких як окуляри MR, які дозволяють користувачам бачити віртуальні об\textquotesingle єкти в реальному світі).

\section{Порівняння AR, VR і MR засвідчує, що усі вони належать до імерсивних технологій, однак між ними наявні принципові відмінності, які зумовлюють різну педагогічну логіку, дидактичний потенціал і специфіку застосування в навчальному процесі. Технологія AR доповнює реальний простір віртуальними елементами, інтегрованими у фізичне середовище, тоді як VR повністю замінює реальний світ штучно створеним цифровим простором, в якому користувач перебуває ізольовано від зовнішніх стимулів. MR поєднує властивості обох попередніх технологій, дозволяючи віртуальним об'єктам не лише відображатися у реальному середовищі, а й взаємодіяти з його елементами в режимі реального часу. AR здебільшого реалізується на доступних мобільних пристроях, таких як смартфони чи планшети, тоді як VR потребує використання спеціалізованих VR-шоломів для забезпечення повного занурення. Натомість MR вимагає застосування високотехнологічних окулярів змішаної реальності, що забезпечують синхронне бачення реального середовища та віртуальних об'єктів, здатних до інтерактивної взаємодії.}\label{ux43fux43eux440ux456ux432ux43dux44fux43dux43dux44f-ar-vr-ux456-mr-ux437ux430ux441ux432ux456ux434ux447ux443ux454-ux449ux43e-ux443ux441ux456-ux432ux43eux43dux438-ux43dux430ux43bux435ux436ux430ux442ux44c-ux434ux43e-ux456ux43cux435ux440ux441ux438ux432ux43dux438ux445-ux442ux435ux445ux43dux43eux43bux43eux433ux456ux439-ux43eux434ux43dux430ux43a-ux43cux456ux436-ux43dux438ux43cux438-ux43dux430ux44fux432ux43dux456-ux43fux440ux438ux43dux446ux438ux43fux43eux432ux456-ux432ux456ux434ux43cux456ux43dux43dux43eux441ux442ux456-ux44fux43aux456-ux437ux443ux43cux43eux432ux43bux44eux44eux442ux44c-ux440ux456ux437ux43dux443-ux43fux435ux434ux430ux433ux43eux433ux456ux447ux43dux443-ux43bux43eux433ux456ux43aux443-ux434ux438ux434ux430ux43aux442ux438ux447ux43dux438ux439-ux43fux43eux442ux435ux43dux446ux456ux430ux43b-ux456-ux441ux43fux435ux446ux438ux444ux456ux43aux443-ux437ux430ux441ux442ux43eux441ux443ux432ux430ux43dux43dux44f-ux432-ux43dux430ux432ux447ux430ux43bux44cux43dux43eux43cux443-ux43fux440ux43eux446ux435ux441ux456.-ux442ux435ux445ux43dux43eux43bux43eux433ux456ux44f-ar-ux434ux43eux43fux43eux432ux43dux44eux454-ux440ux435ux430ux43bux44cux43dux438ux439-ux43fux440ux43eux441ux442ux456ux440-ux432ux456ux440ux442ux443ux430ux43bux44cux43dux438ux43cux438-ux435ux43bux435ux43cux435ux43dux442ux430ux43cux438-ux456ux43dux442ux435ux433ux440ux43eux432ux430ux43dux438ux43cux438-ux443-ux444ux456ux437ux438ux447ux43dux435-ux441ux435ux440ux435ux434ux43eux432ux438ux449ux435-ux442ux43eux434ux456-ux44fux43a-vr-ux43fux43eux432ux43dux456ux441ux442ux44e-ux437ux430ux43cux456ux43dux44eux454-ux440ux435ux430ux43bux44cux43dux438ux439-ux441ux432ux456ux442-ux448ux442ux443ux447ux43dux43e-ux441ux442ux432ux43eux440ux435ux43dux438ux43c-ux446ux438ux444ux440ux43eux432ux438ux43c-ux43fux440ux43eux441ux442ux43eux440ux43eux43c-ux432-ux44fux43aux43eux43cux443-ux43aux43eux440ux438ux441ux442ux443ux432ux430ux447-ux43fux435ux440ux435ux431ux443ux432ux430ux454-ux456ux437ux43eux43bux44cux43eux432ux430ux43dux43e-ux432ux456ux434-ux437ux43eux432ux43dux456ux448ux43dux456ux445-ux441ux442ux438ux43cux443ux43bux456ux432.-mr-ux43fux43eux454ux434ux43dux443ux454-ux432ux43bux430ux441ux442ux438ux432ux43eux441ux442ux456-ux43eux431ux43eux445-ux43fux43eux43fux435ux440ux435ux434ux43dux456ux445-ux442ux435ux445ux43dux43eux43bux43eux433ux456ux439-ux434ux43eux437ux432ux43eux43bux44fux44eux447ux438-ux432ux456ux440ux442ux443ux430ux43bux44cux43dux438ux43c-ux43eux431ux454ux43aux442ux430ux43c-ux43dux435-ux43bux438ux448ux435-ux432ux456ux434ux43eux431ux440ux430ux436ux430ux442ux438ux441ux44f-ux443-ux440ux435ux430ux43bux44cux43dux43eux43cux443-ux441ux435ux440ux435ux434ux43eux432ux438ux449ux456-ux430-ux439-ux432ux437ux430ux454ux43cux43eux434ux456ux44fux442ux438-ux437-ux439ux43eux433ux43e-ux435ux43bux435ux43cux435ux43dux442ux430ux43cux438-ux432-ux440ux435ux436ux438ux43cux456-ux440ux435ux430ux43bux44cux43dux43eux433ux43e-ux447ux430ux441ux443.-ar-ux437ux434ux435ux431ux456ux43bux44cux448ux43eux433ux43e-ux440ux435ux430ux43bux456ux437ux443ux454ux442ux44cux441ux44f-ux43dux430-ux434ux43eux441ux442ux443ux43fux43dux438ux445-ux43cux43eux431ux456ux43bux44cux43dux438ux445-ux43fux440ux438ux441ux442ux440ux43eux44fux445-ux442ux430ux43aux438ux445-ux44fux43a-ux441ux43cux430ux440ux442ux444ux43eux43dux438-ux447ux438-ux43fux43bux430ux43dux448ux435ux442ux438-ux442ux43eux434ux456-ux44fux43a-vr-ux43fux43eux442ux440ux435ux431ux443ux454-ux432ux438ux43aux43eux440ux438ux441ux442ux430ux43dux43dux44f-ux441ux43fux435ux446ux456ux430ux43bux456ux437ux43eux432ux430ux43dux438ux445-vr-ux448ux43eux43bux43eux43cux456ux432-ux434ux43bux44f-ux437ux430ux431ux435ux437ux43fux435ux447ux435ux43dux43dux44f-ux43fux43eux432ux43dux43eux433ux43e-ux437ux430ux43dux443ux440ux435ux43dux43dux44f.-ux43dux430ux442ux43eux43cux456ux441ux442ux44c-mr-ux432ux438ux43cux430ux433ux430ux454-ux437ux430ux441ux442ux43eux441ux443ux432ux430ux43dux43dux44f-ux432ux438ux441ux43eux43aux43eux442ux435ux445ux43dux43eux43bux43eux433ux456ux447ux43dux438ux445-ux43eux43aux443ux43bux44fux440ux456ux432-ux437ux43cux456ux448ux430ux43dux43eux457-ux440ux435ux430ux43bux44cux43dux43eux441ux442ux456-ux449ux43e-ux437ux430ux431ux435ux437ux43fux435ux447ux443ux44eux442ux44c-ux441ux438ux43dux445ux440ux43eux43dux43dux435-ux431ux430ux447ux435ux43dux43dux44f-ux440ux435ux430ux43bux44cux43dux43eux433ux43e-ux441ux435ux440ux435ux434ux43eux432ux438ux449ux430-ux442ux430-ux432ux456ux440ux442ux443ux430ux43bux44cux43dux438ux445-ux43eux431ux454ux43aux442ux456ux432-ux437ux434ux430ux442ux43dux438ux445-ux434ux43e-ux456ux43dux442ux435ux440ux430ux43aux442ux438ux432ux43dux43eux457-ux432ux437ux430ux454ux43cux43eux434ux456ux457.}

Узагальнюючи варто відзначити, що у контексті підготовки майбутніх учителів інформатики імерсивні технології \emph{AR, VR та MR відкривають значний потенціал для модернізації освітнього процесу, розширення можливостей візуалізації, моделювання та створення навчальних ситуацій, недоступних у традиційному середовищі.} Усвідомлення принципових відмінностей між цими технологіями дає змогу обґрунтовано обирати оптимальні інструменти для розробки імерсивних освітніх ресурсів, добудовуючи дидактичні сценарії відповідно до змісту фахової підготовки.

Відтак у межах означеного підходу \emph{віртуальна реальність} характеризується створенням повністю штучного цифрового середовища, у яке користувач занурюється за допомогою спеціалізованих технічних засобів, зокрема VR-шоломів. \emph{Доповнена реальність} ґрунтується на накладанні цифрової інформації на реальний простір, розширюючи можливості його сприйняття та пізнання. Змішана реальність поєднує ознаки віртуальної та доповненої реальності, забезпечуючи інтеграцію цифрових об'єктів у реальний простір та їх взаємодію з фізичним середовищем у режимі реального часу. Окрему групу імерсивних технологій становлять \emph{360° відео й фотографії}, що формують панорамний ефект присутності, а також \emph{інтерактивні симуляції,} спрямовані на моделювання реальних або умовних ситуацій з можливістю прийняття рішень і спостереження за їх наслідками. У сукупності зазначені технології формують імерсивне освітнє середовище, орієнтоване на активізацію пізнавальної діяльності та розширення практичного досвіду здобувачів освіти.

\includegraphics[width=6.69653in,height=4.75717in]{thesis_final/images/media/image3.png}

\emph{Рис. 1.3.} Основні різновиди імерсивних технологій

\emph{\textbf{Імерсивні освітні ресурси}} постають як результат педагогічно вмотивованого застосування імерсивних технологій у освітньому процесі та розглядаються як спеціально спроєктовані цифрові продукти, спрямовані на досягнення визначених освітніх цілей. На відміну від традиційних навчальних матеріалів, такі ресурси ґрунтуються на використанні віртуальної, доповненої, змішаної реальності та інших форм імерсивних середовищ, що забезпечують занурення здобувачів освіти в навчальну діяльність, активну взаємодію з контентом і формування практично орієнтованих умінь у змодельованих або наближених до реальних умовах. Сутнісними характеристиками ІОР є інтерактивність, що передбачає активну участь користувача у взаємодії з навчальним середовищем і вплив на перебіг навчальної діяльності; занурення, яке забезпечує ефект присутності у віртуальному або змішаному освітньому просторі та сприяє глибшому засвоєнню навчального матеріалу; контекстуальність, завдяки якій навчання відбувається у межах реальних або спеціально змодельованих ситуацій, що підвищує практичну значущість набутих знань; візуалізація, орієнтована на використання цифрових візуальних засобів для представлення складних понять і процесів; а також адаптивність, що забезпечує можливість варіативного налаштування змісту й складності ресурсу відповідно до індивідуальних освітніх потреб і рівня підготовки здобувачів освіти.

У межах дослідження було здійснено класифікацію імерсивних освітніх ресурсів з урахуванням їхнього функціонального призначення, дидактичної ролі та ступеня безпосередньої залученості до освітнього процесу. Тобто ІОР класифікуються не за формальною ознакою їхнього типу, а за дидактичною функцією в процесі навчання майбутніх учителів інформатики розробки імерсивних освітніх ресурсів. ІОР вже не лише не як сукупність цифрових продуктів, а більше - інструменти організації навчальної діяльності майбутніх учителів інформатики. Відповідно до запропонованого підходу імерсивні освітні ресурси доцільно поділяти на основні та допоміжні, що відображає їхню різну роль у навчанні розробки імерсивних освітніх ресурсів.

До основних ІОР віднесено імерсивні навчально-методичні комплекси, системи підтримки навчання, навчальні лабораторії та предметні середовища, які безпосередньо забезпечують формування професійних умінь студентів і створюють умови для інтерактивної взаємодії з віртуальними або змішаними середовищами. У межах цієї групи виокремлюються імерсивні програмно-методичні матеріали, засоби оцінювання навчальних досягнень, практикуми, навчально-методичні матеріали та середовища моделювання, що реалізують різні дидактичні функції - від подання навчального контенту до організації симуляційної та проєктної діяльності.

Допоміжні ІОР представлені додатковими науково-навчальними матеріалами та комунікаційними засобами, які розширюють можливості основних ІОР, забезпечують інформаційну підтримку, зворотний зв'язок і координацію навчальної взаємодії. Структуризація ІОР, представлена на Рис.~1.4, спирається на результати попередньої апробації авторського підходу, здійсненої в межах магістерського дослідження \hl{{[}Ш-3{]}} та адаптованої до завдань підготовки майбутніх учителів інформатики.

\includegraphics[width=6.69653in,height=4.70533in]{thesis_final/images/media/image4.png}

\emph{Рис. 1.4.} Функціональна класифікація імерсивних освітніх ресурсів у підготовці майбутніх учителів інформатики

Класифікація відображає послідовність дидактичних функцій імерсивних освітніх ресурсів у процесі навчання майбутніх учителів інформатики, від проєктування та реалізації до рефлексивного оцінювання результатів.

\emph{\textbf{Процес навчання розробки імерсивних освітніх ресурсів}} розглядається як цілісно організована педагогічна діяльність, спрямована на формування у майбутніх учителів інформатики системи знань, умінь і практичних навичок, необхідних для проєктування, створення та впровадження ефективних імерсивних освітніх ресурсів у навчальний процес. Зміст цього процесу ґрунтується на поєднанні \emph{теоретичної підготовки} з \emph{практично орієнтованою діяльністю} та охоплює опанування концептуальних засад імерсивних технологій, зокрема принципів функціонування віртуальної, доповненої та змішаної реальності, а також інших цифрових середовищ занурення. Важливим складником ПНР\_ІОР є \emph{освоєння інструментів розробки} імерсивних ресурсів, що передбачає набуття досвіду роботи з програмними середовищами та платформами для створення VR-, AR- та MR-застосунків, зокрема Unity, Unreal Engine, ARKit, ARCore. Поряд із технічною підготовкою особливе значення надається \emph{педагогічному проєктуванню}, у межах якого майбутні вчителі інформатики навчаються розробляти дидактично обґрунтовані сценарії навчання, визначати освітні цілі та завдання, добирати адекватні методи й прийоми навчальної діяльності з урахуванням можливостей імерсивних середовищ.

Процес навчання також передбачає апробацію створених ІОР шляхом їх тестування та оцінювання, аналізу результатів використання й збирання зворотного зв'язку, що дозволяє коригувати зміст і функціональні характеристики цифрових продуктів. Завершальним компонентом ПНР\_ІОР є розробка навчально-методичного забезпечення, яке охоплює створення інструктивних матеріалів, методичних рекомендацій та супровідної документації для ефективного використання імерсивних освітніх ресурсів у практиці закладів освіти. У сукупності зазначені компоненти забезпечують формування у майбутніх учителів інформатики професійних компетентностей, необхідних для самостійної розробки, педагогічно вмотивованого впровадження та оцінювання імерсивних освітніх ресурсів у сучасному освітньому середовищі.

Отже, на цьому етапі ми припустили можливість структурування моделі процесу навчання розробки імерсивних освітніх ресурсів (Рис. 1.5.)

\includegraphics[width=6.69653in,height=5.77108in]{thesis_final/images/media/image5.png}

\emph{Рис.1.5.} Структурно-функціональна модель процесу навчання розробки імерсивних освітніх ресурсів

\emph{\textbf{Цифровий дизайн}} (далі - ЦД) в освіті та його диференційовані варіанти (digital instructional design, digital learning design, digital media design). Зазначимо, що у вітчизняному науковому дискурсі цифровий дизайн (Digital Design) тлумачиться переважно як ``творчий метод, процес і результат художньо-технічного проектування за допомогою інформаційних технологій промислових виробів та різних комплексів і систем, орієнтованих на найбільшу відповідність створюваних об'єктів потребам сучасної людини та навколишньому природному середовищу як з прагматичної так і з естетичної точки зору'' \hl{{[}З-2, с.11{]}. Таке розуміння зосереджується передусім на візуально-технічному вимірі цифрових продуктів і лише опосередковано враховує їхній педагогічний потенціал.}

\hl{Водночас у науково-педагогічному контексті функціонують близькі, але змістово відмінні варіанти: Зокрема, \emph{Digital instructional design} використовується для позначення проєктування навчального контенту з орієнтацією на результати навчання, структурування змісту та дидактичну доцільність} \hl{{[}{]}; Digital learning design пов'язується з конструюванням цифрового освітнього середовища як цілісної системи навчальних активностей, взаємодії та підтримки здобувачів освіти} \hl{{[}{]}; Digital media design акцентує увагу на візуально-комунікаційному аспекті цифрових ресурсів, мультимедійності та користувацькому досвіді} \hl{{[}{]}. Наявність таких концептуальних розмежувань засвідчує потребу у педагогічному осмисленні цифрового дизайну не лише як техніко-художньої діяльності, а як складника професійної підготовки майбутнього вчителя.}

\hl{Цифровий дизайн у контексті дослідження виконує функцію перехідної ланки між технологічними можливостями імерсивних технологій та їх педагогічно доцільним використанням у навчальному процесі, забезпечуючи трансформацію технічних рішень у дидактично вмотивовані освітні продукти.}

\hl{На перетині цифрового дизайну, педагогічного проєктування та методики навчання інформатики формується \emph{\textbf{професійна готовність майбутнього вчителя інформатики до розробки ІОР}}}\emph{\textbf{.}} Зміст професійної готовності охоплює поєднання теоретичних знань про засади імерсивних технологій, принципи педагогічного проєктування, основи цифрового дизайну та методику навчання інформатики з практичними вміннями проєктувати навчальні сценарії, використовувати інструменти розробки імерсивних ресурсів і створювати інтерактивні навчальні середовища, орієнтовані на активну участь здобувачів освіти. Водночас важливою складовою цієї готовності є особистісні якості майбутнього вчителя інформатики, зокрема креативність, ініціативність, відповідальність, здатність до самостійної діяльності та професійної співпраці, які забезпечують продуктивну роботу в умовах швидкої цифрової трансформації освіти.

Окремого значення набуває мотиваційний аспект професійної готовності, що виявляється у стійкій зацікавленості майбутнього вчителя інформатики у використанні імерсивних технологій в освітньому процесі, прагненні до постійного професійного розвитку та вдосконалення власної педагогічної практики. Саме мотиваційна складова визначає не лише готовність до опанування нових технологій, а й здатність критично осмислювати їх педагогічний потенціал та межі доцільності застосування.

\emph{Професійна готовність майбутнього вчителя інформатики до розробки імерсивних освітніх ресурсів} у межах цього дослідження розглядається як інтегративна характеристика особистості, що виявляється у здатності педагогічно обґрунтовано проєктувати, створювати та впроваджувати імерсивні освітні ресурси на основі поєднання теоретичних знань з імерсивних технологій, цифрового дизайну й методики навчання інформатики, практичних умінь роботи з відповідними інструментами та платформами, а також сформованих особистісно-мотиваційних якостей, необхідних для професійної діяльності в умовах цифрової трансформації освіти. Таке розуміння професійної готовності узгоджується з положеннями стандарту підготовки вчителя інформатики, відповідно до яких фахова діяльність педагога передбачає не лише володіння цифровими технологіями, а й здатність до їх педагогічно доцільного використання для досягнення результатів навчання, проєктування освітнього середовища та забезпечення активної пізнавальної діяльності здобувачів освіти. Подальший розвиток AR, VR і MR зумовлює зростання їхнього значення у формуванні цифрової та проєктувальної компетентності студентів, а також посилює інтеграцію імерсивних рішень у професійну діяльність сучасного вчителя. У цьому контексті цифровий дизайн імерсивних середовищ постає як складник професійної діяльності вчителя інформатики, що поєднує технологічну, дидактичну та методичну компетентності.

\includegraphics[width=6.69653in,height=2.5455in]{thesis_final/images/media/image6.png}

\emph{Рис. 1.6.} Структурно-логічна схема професійної готовності майбутнього вчителя інформатики до розробки імерсивних освітніх ресурсів

У сукупності зазначені на Рис. 1.6. компоненти професійної готовності утворюють підґрунтя для переходу до осмислення педагогічної ефективності імерсивних освітніх ресурсів, оскільки рівень такої ефективності безпосередньо залежить від здатності вчителя інтегрувати цифровий дизайн імерсивних середовищ у логіку навчального процесу та орієнтувати їх використання на досягнення визначених освітніх результатів.

\emph{Педагогічна ефективність імерсивних освітніх ресурсів} у межах цього дослідження розглядається як міра їх здатності забезпечувати досягнення визначених освітніх результатів за умови педагогічно доцільного проєктування та використання в навчальному процесі. На відміну від суто технологічної результативності, педагогічна ефективність ІОР зумовлюється не стільки рівнем їх технічної складності, скільки відповідністю дидактичним цілям, змісту навчання та особливостям пізнавальної діяльності здобувачів освіти.

У педагогічному вимірі ІОР набувають педагогічної цінності лише за умови їх інтеграції у логіку навчального процесу, що передбачає чітке визначення освітніх цілей, продумане сценарне проєктування навчальної діяльності, адекватний добір методів і форм роботи, а також рефлексивне оцінювання результатів навчання. Відтак педагогічна ефективність ІОР безпосередньо залежить від рівня професійної готовності вчителя інформатики, який виступає суб'єктом їх розробки, адаптації та впровадження.

Педагогічно ефективні ІОР сприяють активізації пізнавальної діяльності, формуванню стійкого навчального інтересу, розвитку практичних умінь і навичок, а також забезпечують контекстуальність навчання через моделювання реальних або наближених до реальних ситуацій. Водночас відсутність педагогічного проєктування або формальне використання ІТ зводить їх потенціал до рівня візуального доповнення навчального матеріалу, не забезпечуючи суттєвого впливу на якість освітніх результатів.

Отже, педагогічна ефективність імерсивних освітніх ресурсів постає не автономною характеристикою цифрового продукту, а похідною від цілісного процесу навчання їх розробки та впровадження. Саме у взаємозв'язку з професійною готовністю майбутнього вчителя інформатики вона розглядається як важливий критерій обґрунтування доцільності використання імерсивних технологій у сучасному освітньому середовищі.

Можемо констатувати, що педагогічна ефективність імерсивних освітніх ресурсів -- це ступінь досягнення освітніх цілей за допомогою використання ІОР. Вона може бути оцінена за різними критеріями: \emph{засвоєння знань} (рівень розуміння та запам\textquotesingle ятовування навчального матеріалу), \emph{формування навичок} (розвиток практичних вмінь та навичок, необхідних для вирішення професійних завдань); \emph{мотивація до навчання} (підвищення інтересу до навчання та активності учнів); \emph{залучення} (творення ефекту присутності та залучення учнів до навчального процесу); \emph{зручність використання} (простота та зручність використання ІОР), що сприяє ефективному навчанню. Оцінка педагогічної ефективності ІОР є важливим етапом їх розробки та впровадження в освітній процес.

Водночас маємо зазначити: аналіз сучасних наукових джерел засвідчує (Розділ~1.1.), що використання імерсивних технологій в освітньому процесі характеризується поєднанням виразного педагогічного потенціалу та низки суперечливих наслідків, які зумовлюють необхідність їх педагогічно виваженого впровадження. Дослідники наголошують, що імерсивні технології здатні істотно підвищувати якість навчання на різних рівнях освіти, насамперед за рахунок створення ефекту присутності та занурення, активізації пізнавальної діяльності й зростання навчальної мотивації \hl{{[}Г-1, с. 138--139{]}.} Їх привабливість для здобувачів освіти пов'язується з можливістю інтерактивної взаємодії з навчальним контентом, розвитку творчого мислення та формування ключових компетентностей \hl{{[}А-1, с. 58{]}.}

Окремі дослідження підкреслюють, що імерсивне навчання трансформує традиційні підходи до засвоєння знань, перетворюючи навчальний процес на інтуїтивно зрозумілу, контекстуалізовану діяльність, у межах якої здобувачі освіти можуть експериментувати, здійснювати практичні дії та закріплювати знання у змодельованих ситуаціях \hl{{[}Л-9, с. 55--56{]}.} Значущим аспектом позитивного впливу ІТ вважається їх потенціал у зниженні психологічної напруги, характерної для онлайн-навчання, а також у підтримці навчальної успішності в умовах цифровізації освітнього середовища \hl{{[}Л-9, с. 56{]}.}

Водночас системний аналіз літератури дозволяє виявити низку ризиків і проблем, що супроводжують упровадження імерсивних технологій. Серед них науковці виокремлюють неготовність здобувачів освіти, батьків і педагогів до педагогічно обґрунтованого використання комп'ютерно орієнтованих та імерсивних навчальних систем, що знижує ефективність їх застосування та може призводити до формалізації інновацій \hl{{[}Г-1, с.138--139{]}.} Також наголошується на високій вартості обладнання й програмного забезпечення, складності технічного обслуговування та ризиках надмірної віртуалізації освітнього середовища, що потенційно обмежує живу комунікацію та соціальну взаємодію учасників освітнього процесу \hl{{[}Г-1, с. 138{]}.}

\hl{Зарубіжні вчені констатують існуванні двох протилежних позицій у науковому дискурсі: з одного боку, припущення про позитивний вплив IVR через підвищення залученості та емоцій, з іншого - гіпотеза про негативний вплив через когнітивне перевантаження та відволікання {[}Ре-2{]}. Отже} окрему групу суперечливих наслідків становлять питання впливу імерсивних технологій на психофізичний стан здобувачів освіти, зокрема можливі негативні ефекти тривалого перебування у віртуальних середовищах, що розглядаються дослідниками як аспект, який потребує подальшого наукового вивчення та нормативного врегулювання \hl{{[}Г-1, с. 138{]}.}

Узагальнено позиції вчених що впливу ІТ на освітній процес подано у Таблиці 1.1.

\emph{Таблиця 1.1.}

\textbf{Позитивні та суперечливі наслідки використання ІТ в освіті}

{\def\LTcaptype{none} % do not increment counter
\begin{longtable}[]{@{}
  >{\centering\arraybackslash}p{(\linewidth - 4\tabcolsep) * \real{0.2357}}
  >{\centering\arraybackslash}p{(\linewidth - 4\tabcolsep) * \real{0.4096}}
  >{\centering\arraybackslash}p{(\linewidth - 4\tabcolsep) * \real{0.3383}}@{}}
\toprule\noalign{}
\endhead
\bottomrule\noalign{}
\endlastfoot
\textbf{Аспект} & \textbf{Позитивні наслідки} & \textbf{Суперечливі наслідки / ризики} \\
Дидактичний & Поглиблене засвоєння матеріалу; контекстуальність навчання; відпрацювання практичних навичок у безпечному середовищі {[}Г-1; Л-9{]} & Формалізація використання без педагогічного проєктування; перевантаження візуальними ефектами \\
Мотиваційний & Зростання навчальної мотивації; зацікавленість; розвиток творчого мислення {[}А-1; Л-9{]} & Залежність мотивації від технологічної новизни \\
Особистісний & Розвиток соціальних навичок, співпраці, комунікації у віртуальних середовищах {[}Л-9; Thompson, 2021{]} & Зниження рівня безпосередньої соціальної взаємодії \\
Організаційний & Розширення доступу до ресурсів; подолання географічних і матеріальних обмежень {[}Г-1{]} & Висока вартість обладнання; складність технічного супроводу \\
Професійно-

педагогічний & Можливості модернізації підготовки майбутніх учителів; розвиток цифрових і педагогічних компетентностей {[}А-1{]} & Неготовність педагогів; потреба у спеціальній підготовці \\
Здоров'я-

збережувальний & Потенціал зниження психологічної напруги в онлайн-навчанні {[}Л-9{]} & Неоднозначний вплив на психофізичний стан здобувачів освіти {[}Г-1{]} \\
\end{longtable}
}

Подані у таблиці відомості засвідчують, що педагогічний потенціал імерсивних технологій реалізується у різних вимірах освітнього процесу -- дидактичному, мотиваційному, особистісному, організаційному та професійно-педагогічному. Водночас кожен із позитивних ефектів супроводжується певними ризиками, що зумовлює необхідність системного та педагогічно виваженого підходу до впровадження імерсивних технологій в освіті. Така подвійність результатів підкреслює, що ефективність використання імерсивних технологій визначається не стільки рівнем їх технологічної складності, скільки умовами педагогічного проєктування, підготовленістю вчителя та методичною доцільністю інтеграції у навчальний процес.

Узагальнення результатів теоретичного аналізу дозволяє представити імерсивні технології в освіті як цілісне педагогічне явище, що поєднує технологічні, дидактичні та особистісно-професійні компоненти (Рис.1.5. Структурно-логічна схема імерсивних технологій, представлених в освіті). У такому взаємозв'язку імерсивні технології постають не лише як сукупність цифрових інструментів, а як основа для проєктування імерсивних освітніх ресурсів, організації процесу навчання їх розробки, формування професійної готовності майбутнього вчителя інформатики та досягнення педагогічної ефективності навчання.

\includegraphics[width=6.69653in,height=3.623in]{thesis_final/images/media/image7.png}

\emph{Рис}.\emph{~1.7.} Структурно-логічна схема імерсивних технологій, представлених в освіті.

Представлена структурно-логічна схема відображає взаємозв'язок технологічних, дидактичних і особистісно-професійних компонентів, що визначають потенціал імерсивних технологій у сучасному освітньому середовищі. Разом із тим теоретичне осмислення імерсивних технологій потребує доповнення аналізом практичних способів їх реалізації в освітньому процесі. Саме методичні підходи до використання імерсивних технологій визначають, яким чином технологічні можливості трансформуються у педагогічні рішення, що забезпечують досягнення результатів навчання, урахування індивідуальних освітніх потреб та підвищення якості освітнього процесу. З огляду на це доцільним є звернення до вітчизняного досвіду, у межах якого напрацьовано різноманітні методичні підходи до інтеграції імерсивних технологій у навчання. Їх вивчення дозволяє виявити усталені практики, інноваційні рішення та тенденції, що мають значення для подальшого обґрунтування методичних засад використання імерсивних технологій у вітчизняній освітній практиці, чому й присвячено наступний підрозділ.

\textbf{1.3. Вітчизняний досвід використання імерсивних освітніх ресурсів у професійній підготовці майбутніх учителів інформатики}

З метою вивчення та узагальнення наявного вітчизняного досвіду використання імерсивних освітніх ресурсів у підготовці майбутніх учителів інформатики було здійснено аналіз офіційних сайтів провідних закладів вищої освіти України, у яких реалізуються освітні програми з підготовки фахівців відповідного профілю. Увагу зосереджено на змісті освітньо-професійних програм, навчальних планах, описах освітніх компонентів, матеріалах навчально-методичного забезпечення, а також на відкритих інформаційних ресурсах, що відображають практики використання цифрових та імерсивних технологій у професійній підготовці майбутніх учителів інформатики.

Передусім звертаємо увагу на відсутність чинного державного стандарту вищої освіти першого (бакалаврського) рівня за спеціальністю 014 «Середня освіта (за предметними спеціальностями)», що зумовлює необхідність звернення до проєкту відповідного стандарту як нормативно-орієнтаційної основи дослідження \hl{{[}{]}}. Аналіз положень проєкту стандарту вищої освіти України засвідчує його спрямованість на формування компетентнісної моделі підготовки майбутнього вчителя, у межах якої цифрові технології розглядаються як невід'ємний складник професійної діяльності.

У переліку загальних і фахових компетентностей проєкту стандарту акцент зроблено на здатності застосовувати інформаційно-комунікаційні технології в освітній діяльності, \hl{добирати і використовувати сучасні та ефективні методики і технології навчання, виховання і розвитку учнів; використовувати цифрові технології в освітньому процесі, створювати та використовувати цифрові освітні ресурси \textbf{(ЗК17; ФК6; ФК11)} {[}{]}\textbf{.}} Наведені формулювання корелюють із функціональними можливостями ІОР, зокрема у частині візуалізації складних об'єктів і процесів, моделювання навчальних ситуацій, організації практикоорієнтованої діяльності. Водночас у проєкті стандарту відсутні прямі посилання на імерсивні технології або відповідні цифрові середовища, що свідчить про їх опосередковану, а не цілеспрямовано визначену присутність у нормативному полі підготовки майбутніх учителів.

Аналіз очікуваних результатів навчання дозволяє констатувати, що проєкт закладає вимоги до сформованості вмінь використовувати цифрові інструменти для організації навчального процесу, оцінювання навчальних досягнень і професійного саморозвитку \hl{(\textbf{РН15).}} Результати навчання створюють можливість для включення імерсивних освітніх ресурсів як засобів реалізації відповідних педагогічних завдань, однак не визначають конкретних шляхів, форм або рівнів їх застосування.

Проєкт Стандарту вибудувано як універсальну рамку, що охоплює весь спектр спеціалізацій у межах відповідної галузі підготовки, водночас його зміст не обмежується лише загальними орієнтирами. Документ містить конкретизований перелік програмних результатів навчання, диференційований за окремими спеціалізаціями, зокрема для 014.09 «Середня освіта (Інформатика)», що свідчить про поєднання рамкового підходу зі спробою врахування специфіки фахової підготовки майбутніх учителів інформатики.

Поданий у Стандарті перелік ПРН репрезентує переважно інструментально-технологічну логіку підготовки майбутнього вчителя інформатики. Зміст ПРН орієнтований на сформованість операційних умінь роботи з даними (ПРН1), програмним забезпеченням і цифровими продуктами (ПРН2), комп'ютерною технікою та мережевою інфраструктурою (ПРН3), а також на моделювання й проведення комп'ютерного експерименту (ПРН4), що відображає класичну модель фахової ІТ-підготовки в межах спеціалізації 014.09 «Середня освіта (Інформатика)» (А4.09 відповідно до оновлення Переліку галузей знань і спеціальностей, затвердженого постановою Кабінету Міністрів України від 30 серпня 2024 р. №1021 «Про внесення змін до переліку галузей знань і спеціальностей, за якими здійснюється підготовка здобувачів вищої та фахової передвищої освіти'' \hl{{[}{]}).} Водночас ці ПРН лише опосередковано пов'язані з педагогічним контекстом та майже не артикулюють освітній вимір використання цифрових технологій. Артикульовано створення цифрових продуктів (ПРН2), розв'язування задач з інформатики різного рівня складності (ПРН5), моделювання (ПРН4) та забезпечення інформаційної безпеки (ПРН6), однак відсутнє розмежування між використанням цифрових засобів як об'єкта вивчення і як дидактичного інструмента організації навчального процесу.

У контексті теми навчання розробки імерсивних освітніх ресурсів ці ПРН створюють лише базове технологічне підґрунтя. Вони формують передумови для опанування інструментів програмування \textbf{(ПРН2)}, інформаційного моделювання та комп'ютерного експерименту \textbf{(ПРН4)}, які потенційно можуть бути використані під час створення VR/AR-рішень, симуляцій або інтерактивних навчальних середовищ. Проте сам процес педагогічно вмотивованого проєктування імерсивних освітніх ресурсів, їх дидактичний дизайн і орієнтація на досягнення результатів навчання здобувачів освіти у жодному з ПРН прямо не зафіксовані. Отже, перелік ПРН засвідчує технологічну готовність випускника до роботи з цифровими системами (ПРН1--ПРН6), але не відображає сформованості вмінь педагогічного проєктування, методично обґрунтованого використання імерсивних технологій та інтеграції таких ресурсів у структуру шкільного курсу інформатики. Окреслене актуалізує потребу у доповненні або уточненні ПРН з урахуванням освітньо-дидактичної специфіки імерсивних освітніх ресурсів та ролі вчителя інформатики як їх розробника і впроваджувача.

Відповідно питання системності й методичної доцільності використання імерсивних ресурсів залишаються на розсуд ЗВО та розробників освітньо-професійних програм (далі - ОПП). Проєкт стандарту виконує функцію нормативної рамки, що задає загальні орієнтири професійної підготовки майбутніх учителів інформатики, але не містить спеціалізованих вимог щодо використання імерсивних освітніх ресурсів. Зазначене зумовлює варіативність підходів до їх упровадження в ОПП і робочих програмах (далі - РП) навчальних дисциплін та актуалізує необхідність науково аналізу й узагальнення наявного вітчизняного досвіду використання імерсивних ресурсів у підготовці майбутніх учителів інформатики.

Задля з'ясування реального стану та особливостей вітчизняного досвіду підготовки майбутніх учителів інформатики було використано інформаційні можливості \textbf{Єдиної державної електронної бази з питань освіти} (далі - ЄДЕБО; https://info.edbo.gov.ua/), зокрема розділ пошуку конкурсних пропозицій (Вступ\_2025), що дозволяє ідентифікувати заклади вищої освіти України, у яких здійснюється підготовка здобувачів за спеціальністю А4.09 «Середня освіта (Інформатика)», а також проаналізувати актуальні освітні програми й умови їх реалізації (Рис. 1.8).

\includegraphics[width=5.44119in,height=3.39521in]{thesis_final/images/media/image8.png}

\emph{Рис. 1.8.} Інтерфейс пошуку конкурсних пропозицій у Єдиній державній електронній базі з питань освіти (ЄДЕБО)

За результатами опрацювання інформації, представленої в ЄДЕБО, встановлено, що підготовка майбутніх учителів інформатики здійснюється у 38 закладах вищої освіти України, у межах яких запропоновано 106 конкурсних пропозицій за відповідною спеціальністю. Така кількість освітніх пропозицій свідчить про широку інституційну представленість підготовки майбутніх учителів інформатики у вітчизняному освітньому просторі та створює підґрунтя для виявлення різних підходів до використання імерсивних освітніх ресурсів у професійній підготовці.

Аналіз переліку закладів вищої освіти, у яких реалізуються відповідні конкурсні пропозиції, засвідчує, що підготовка майбутніх учителів інформатики здійснюється у 38 закладах вищої освіти, серед яких переважають педагогічні університети (18 закладів), поряд із класичними університетами (14 закладів) та багатопрофільними й спеціалізованими ЗВО (6 закладів). Представленість ОПП у різних регіонах України вказує на територіальну розгалуженість підготовки фахівців зазначеного профілю та водночас актуалізує питання неоднорідності умов упровадження інноваційних освітніх практик, зокрема імерсивних освітніх ресурсів. Припускаємо, що така інституційна різнорідність зумовлює відмінності в освітній логіці, ресурсному забезпеченні та методичних підходах до організації освітнього процесу. За таких обставин практики їх використання формуються нерівномірно та не завжди набувають системного характеру в межах освітніх програм.

З метою встановлення особливостей наявного вітчизняного досвіду використання імерсивних освітніх ресурсів у підготовці майбутніх учителів інформатики подальше дослідження було зорієнтоване на аналіз змісту освітніх програм підготовки бакалаврів за спеціальністю 014.09 «Середня освіта (Інформатика)», а також робочих програм навчальних дисциплін, у межах яких реалізується професійна підготовка здобувачів освіти. Обраний підхід дозволив виявити не лише декларативну присутність цифрових та імерсивних технологій у програмних документах, а й конкретні форми, напрями та педагогічні функції їх використання у освітньому процесі.

Аналіз ОПП і робочих програм навчальних дисциплін здійснювався з урахуванням змістових і структурних критеріїв, що дозволили виявити реальний стан використання ІОР у підготовці майбутніх учителів інформатики. Увагу було зосереджено на наявності в програмних документах положень, які відображають застосування ІОР у формулюваннях результатів навчання, переліках програмних компетентностей, описах змісту освітніх компонентів і видах навчальної діяльності здобувачів освіти. Окремо враховувалися відомості про використання цифрових середовищ моделювання, віртуальних лабораторій, симуляційних і проєктних форм роботи, а також особливості організації практичної підготовки.

Перед презентацією результатів аналізу доцільно зауважити, що на офіційних вебсторінках окремих українських закладів вищої освіти, зокрема Класичний приватний університет, Мелітопольський державний педагогічний університет імені Богдана Хмельницького, Національний університет «Чернігівський колегіум» імені Т. Г. Шевченка та Сумський державний педагогічний університет імені А. С. Макаренка, - відсутнє оприлюднення ОПП за відповідною спеціальністю, що ускладнило повноцінний порівняльний розгляд програмного забезпечення підготовки фахівців.

Порівняльний аналіз ОПП підготовки майбутніх учителів інформатики, здійснений на матеріалі програм закладів вищої освіти України, засвідчив наявність стійкої тенденції до декларативного визнання ролі цифрових технологій у професійній підготовці за відсутності їх цілісної концептуалізації в площині імерсивної освіти. Практично в усіх проаналізованих програмах цифрові технології розглядаються як універсальний інструмент навчання, що забезпечує модернізацію освітнього процесу, проте не як окремий предмет педагогічного проєктування чи спеціально організований напрям професійної підготовки. Так, у програмах педагогічних університетів, зокрема Бердянський державний педагогічний університет, Житомирський державний університет імені Івана Франка, Херсонський державний університет, використання цифрових технологій і комп'ютерного моделювання декларується на рівні загальних результатів навчання та фахових компетентностей, однак імерсивні освітні ресурси не виділяються як самостійний об'єкт опанування або як специфічний результат професійної підготовки.

Така ж ситуація простежується і в програмах, що мають підвищений інноваційний або STEM-орієнтований потенціал. У низці ОПП, зокрема Волинський національний університет імені Лесі Українки, Черкаський національний університет імені Богдана Хмельницького, Криворізький державний педагогічний університет, зафіксовано наявність дисциплін, пов'язаних із комп'ютерною графікою, цифровим дизайном, STEM/STEAM-освітою, робототехнікою та комп'ютерним моделюванням, що створює сприятливі технологічні передумови для інтеграції імерсивних освітніх ресурсів. Водночас навіть за таких умов імерсивність залишається імпліцитною характеристикою цифрового навчального середовища, а не концептуально оформленим елементом освітньої програми.

Показовим є й те, що в жодній з проаналізованих ОПП імерсивні освітні ресурси не зафіксовані як самостійна мета підготовки, окремий блок результатів навчання або структурний компонент освітньої траєкторії студента. Навіть у програмах, де наявні курси з комп'ютерної анімації, доповненої реальності чи мультимедійних технологій, як, наприклад, у Кременчуцький національний університет імені Михайла Остроградського, імерсивні засоби розглядаються переважно як окремі технологічні інструменти, а не як педагогічно осмислений ресурс, що потребує спеціальної методики розробки та застосування в освітньому процесі.

Відсутність у проаналізованих ОПП чіткої фіксації імерсивних освітніх ресурсів як самостійного об'єкта професійної підготовки, на нашу думку, свідчить про початковий етап інституційного осмислення імерсивності в педагогічній освіті. Наявні програми акумулюють окремі цифрові, технологічні й змістові елементи, що потенційно можуть слугувати підґрунтям для розвитку імерсивної освіти, однак не трансформують їх у цілісну педагогічну модель підготовки майбутнього вчителя інформатики.

Водночас аналіз формулювань \textbf{результатів навчання} демонструє переважно імпліцитний характер присутності ідей імерсивності, що проявляється у широкому, узагальненому трактуванні цифрових технологій як інструментів професійної діяльності майбутнього вчителя інформатики. Переважна більшість ОПП орієнтована на формування здатності застосовувати інформаційно-комунікаційні технології, цифрові освітні ресурси, засоби комп'ютерного моделювання, візуалізації та програмування, не конкретизуючи характер, рівень і педагогічне призначення таких технологій у навчальному процесі. Зокрема, у програмах Б\textbf{е}рдянського державного педагогічного університету, Ніжинського державного університету імені Миколи Гоголя, Рівненського державного гуманітарного університету результати навчання формулюються через здатність «використовувати сучасні ІКТ», «застосовувати цифрові освітні ресурси» або «реалізовувати комп'ютерне моделювання в освітньому процесі», що створює технологічне підґрунтя для імерсивних практик, однак не окреслює їх як окремий педагогічний результат.

Подібна логіка простежується і в програмах з підвищеним інноваційним або STEM-орієнтованим спрямуванням. Так, у формулюваннях результатів навчання ОПП Волинського національного університету імені Лесі Українки, Черкаського національного університету імені Богдана Хмельницького, Криворізького державного педагогічного університету акцент робиться на здатності здійснювати цифрове та STEM-моделювання, проєктувати освітні продукти, застосовувати мультимедійні та інноваційні технології навчання. Водночас ІОР не фіксуються як специфічний об'єкт професійної діяльності, а розглядаються лише як можливий похідний інструмент у межах загального цифрового середовища.

Показовим є й те, що навіть у програмах, де результати навчання пов'язані з розробленням складних цифрових продуктів, комп'ютерної графіки або візуалізації навчального контенту, (Кременчуцький національний університет імені Михайла Остроградського чи Український державний університет імені Михайла Драгоманова) імерсивність не набуває статусу окремо визначеного результату навчання. Формулювання ПРН зосереджуються на загальній здатності застосовувати цифрові технології та інноваційні засоби навчання, залишаючи поза увагою питання педагогічного проєктування, методичного супроводу й оцінювання ефективності використання імерсивних освітніх ресурсів.

Отже, результати навчання, задекларовані в більшості проаналізованих ОПП, створюють функціональні передумови для використання імерсивних освітніх ресурсів, однак не забезпечують їх педагогічної, дидактичної та методичної конкретизації як окремої складової професійної компетентності майбутнього вчителя інформатики. Така ситуація зумовлює розрив між технологічним потенціалом цифрового навчального середовища та рівнем його цілеспрямованого педагогічного осмислення в системі фахової підготовки.

Аналіз змісту навчальних дисциплін дозволяє виявити фрагментарну та потенційну інтеграцію імерсивних освітніх ресурсів, зосереджену переважно в курсах технологічного та інноваційного спрямування. Найчастіше відповідний потенціал акумулюється в дисциплінах з комп'ютерної графіки, комп'ютерного моделювання, мультимедійних технологій, STEM/STEAM-освіти, робототехніки, вебтехнологій, а також у вибіркових освітніх компонентах, зорієнтованих на цифровий дизайн і розроблення освітніх продуктів. Так, у низці ЗВО, зокрема Волинський національний університет імені Лесі Українки та Житомирський державний університет імені Івана Франка, потенціал використання імерсивних освітніх ресурсів закладено в дисциплінах «Комп'ютерна графіка», «Мультимедійні технології», «Комп'ютерне моделювання», де передбачено роботу з тривимірними об'єктами, візуалізацією процесів і цифровим моделюванням навчального контенту.

У програмах із виразнішим STEM-орієнтованим компонентом, зокрема Черкаський національний університет імені Богдана Хмельницького, Криворізький державний педагогічний університет, Кременецька обласна гуманітарно-педагогічна академія імені Тараса Шевченка, відповідні можливості простежуються в межах дисциплін «STEM-освіта», «STEM/STEAM-технології», «Освітня робототехніка», «3D-моделювання», «Основи комп'ютерної анімації». Зміст цих курсів передбачає створення цифрових моделей, симуляцій і віртуальних об'єктів, що технологічно наближує здобувачів освіти до роботи з імерсивними середовищами, однак не супроводжується педагогічним осмисленням специфіки імерсивного навчання як окремого дидактичного феномена.

Показовими є також приклади вибіркових дисциплін інноваційного спрямування. У програмах Кременчуцький національний університет імені Михайла Остроградського та Український державний університет імені Михайла Драгоманова зафіксовано курси «Комп'ютерна графіка та доповнена реальність», «Інноваційні технології навчання інформатики», «Освітня робототехніка», «Цифрові освітні технології», у межах яких можливе використання елементів доповненої або віртуальної реальності. Водночас аналіз програм дисциплін засвідчує, що такі елементи подаються як окремі технологічні інструменти або ілюстративні засоби, а не як об'єкт цілеспрямованого педагогічного проєктування. Отже, навіть за наявності дисциплін із підвищеним технологічним потенціалом інтеграція імерсивних освітніх ресурсів у зміст підготовки майбутніх учителів інформатики має епізодичний і несистемний характер. Вона не супроводжується чітко визначеними дидактичними цілями, методами навчання, логікою поетапного формування відповідних умінь або критеріями оцінювання результатів, що унеможливлює перехід від фрагментарного використання імерсивних технологій до цілісної моделі їх педагогічного застосування.

\emph{\textbf{Форми використання цифрових та потенційно імерсивних ресурсів}} у проаналізованих ОПП зводяться переважно до інструментального, демонстраційного або проєктного рівня, що обмежує їх педагогічний потенціал рамками допоміжних засобів навчання. У більшості програм імерсивні ресурси залучаються в межах лабораторних і практичних занять з комп'ютерної графіки, комп'ютерного моделювання, мультимедійних технологій і програмування, де вони використовуються для візуалізації об'єктів, ілюстрації процесів або відпрацювання окремих технічних операцій. Такий підхід простежується, зокрема, в освітніх програмах Бердянського державного педагогічного університету, Рівненського державного гуманітарного університету, Херсонського державного університету, де цифрові засоби застосовуються насамперед як інструменти підтримки засвоєння навчального матеріалу, а не як самостійний об'єкт педагогічного аналізу.

Проєктний рівень використання цифрових ресурсів представлений переважно у формі курсових робіт, індивідуальних або групових проєктів, STEM- і STEAM-завдань, у межах яких здобувачі освіти розробляють електронні навчальні матеріали, цифрові моделі, інтерактивні презентації або елементи навчальних середовищ. Подібні форми зафіксовано в програмах Волинського національного університету імені Лесі Українки, Черкаського національного університету імені Богдана Хмельницького, Криворізького державного педагогічного університету, де цифрові та візуалізаційні засоби використовуються в процесі проєктування навчального контенту. Водночас навіть у таких випадках ІОР не розглядаються як специфічний дидактичний інструмент, що потребує окремої методики навчання, а сприймаються як похідний продукт загальної цифрової діяльності.

Використання цифрових ресурсів під час педагогічних практик також має переважно демонстраційний або допоміжний характер. Майбутні вчителі застосовують мультимедійні презентації, відеофрагменти, інтерактивні вправи або цифрові симуляції для пояснення навчального матеріалу, однак такі практики не передбачають системного залучення імерсивних середовищ як засобу організації навчальної взаємодії чи формування досвіду навчання в імерсивному форматі. Подібна логіка простежується, зокрема, у програмах Ніжинського державного університету імені Миколи Гоголя, Уманського державного педагогічного університету імені Павла Тичини, Житомирського державного університету імені Івана Франка.

У проаналізованих ОПП імерсивні освітні ресурси фактично не набувають статусу самостійного дидактичного інструмента, здатного трансформувати логіку навчального процесу, структуру професійної діяльності майбутнього вчителя інформатики та способи організації освітньої взаємодії. Їх використання обмежується рівнем технічної підтримки навчання, що актуалізує потребу в розробленні спеціально організованої методики навчання, зорієнтованої на цілеспрямоване формування вмінь проєктування й педагогічного застосування імерсивних освітніх ресурсів.

Результати порівняльного аналізу дають підстави стверджувати, що сучасні освітньо-професійні програми підготовки майбутніх учителів інформатики перебувають на етапі накопичення технологічних і змістових передумов для впровадження імерсивних освітніх ресурсів, однак ще не демонструють їх системної педагогічної інтерпретації. Виявлений розрив між наявністю цифрово-інструментального потенціалу та відсутністю цілісної методики навчання розробки імерсивних освітніх ресурсів актуалізує необхідність обґрунтування спеціально організованого педагогічного процесу, що передбачає цілеспрямоване формування відповідних умінь, компетентностей і професійних дій майбутніх учителів інформатики.

Узагальнені результати порівняльного аналізу освітньо-професійних програм представлено в зведеній таблиці, яка подається у \hl{Додатку А} як аналітичне узагальнення.

Результати аналізу дають підстави констатувати наявність системного розриву між технологічним потенціалом освітніх програм і рівнем його педагогічної реалізації. Сучасні освітньо-професійні програми акумулюють окремі елементи цифрової та імерсивно орієнтованої діяльності, однак не забезпечують цілісної моделі навчання розробки й педагогічного застосування імерсивних освітніх ресурсів. Вітчизняний досвід їх застосування найбільш виразно простежується у системі підготовки майбутніх учителів інформатики, де імерсивні ресурси поступово інтегруються в освітні програми як засоби формування проєктно-конструкторських, технологічних і рефлексивних умінь. Водночас практика їх використання залишається нерівномірною та переважно фрагментарною, що актуалізує потребу в узагальненні й науково-педагогічному осмисленні наявних підходів.

У сучасних українських освітніх програмах з підготовки майбутніх учителів інформатики імерсивні освітні ресурси застосовуються як засоби візуалізації складних процесів, моделювання навчальних і професійних ситуацій, організації практикоорієнтованої діяльності та підтримки взаємодії суб'єктів освітнього процесу. Водночас їх використання має переважно фрагментарний характер і відзначається різним рівнем методичної обґрунтованості, що актуалізує потребу в узагальненні наявного вітчизняного досвіду та його науково-педагогічному осмисленні.

Виявлені суперечності обґрунтовують необхідність розроблення спеціально організованої методики навчання, спрямованої на формування в майбутніх учителів інформатики здатності цілеспрямовано проєктувати, дидактично осмислювати та педагогічно використовувати імерсивні освітні ресурси в освітньому процесі.

\textbf{Висновки до першого розділу}

У розділі окреслено концептуальні засади навчання майбутніх учителів інформатики розробки імерсивних освітніх ресурсів як перспективного напряму педагогічної інноватики в умовах цифрової трансформації освіти. Обґрунтовано, що створення та впровадження VR/AR/XR-рішень у навчальний процес виходить за межі суто технічної діяльності й потребує інтеграції технологічних, дидактичних і дизайнерських складників у структурі професійної підготовки вчителя.

Узагальнення результатів аналізу наукових джерел за п'ятьма проблемно-тематичними кластерами дозволяє стверджувати, що імерсивні освітні ресурси формують цілісне міждисциплінарне поле досліджень, у межах якого поєднуються загальнотеоретичні та дидактичні засади, освітні практики різних рівнів, питання професійної педагогічної підготовки, інклюзивний вимір і галузево-порівняльні підходи. Така структуризація наукового доробку засвідчує еволюцію імерсивних технологій від інноваційного інструменту навчання до системного чинника трансформації освітнього середовища, педагогічної діяльності та освітньої доступності, що створює методологічне підґрунтя для подальшого теоретичного обґрунтування й розроблення авторських підходів у межах досліджуваної проблеми.

Теоретичне узагальнення наукових джерел дозволяє розглядати імерсивні технології як нову дидактичну реальність, здатну забезпечити багаторівневу взаємодію здобувачів освіти з навчальним змістом, активізувати когнітивні процеси та посилити практико-орієнтовану спрямованість професійної підготовки.

У межах сформованого теоретичного поля окреслено структурно-функціональні ознаки імерсивності, прийоми інтерактивної візуалізації, особливості ефектів занурення та присутності, а також педагогічні ефекти їх застосування у процесі опанування інформатичних і міждисциплінарних знань. Імерсивні освітні ресурси постають як органічна складова цифрової трансформації освіти, що ґрунтується на принципах інтерактивності, занурення та візуалізації навчального контенту й зумовлює переосмислення змісту, форм і засобів професійної підготовки майбутнього вчителя.

Уточнення сутнісних характеристик імерсивних освітніх ресурсів у межах дослідження дало змогу розглядати їх не як окремі цифрові продукти або технічні інновації, а як інтерактивні дидактичні засоби нового покоління, побудовані на моделюванні віртуальних, доповнених і змішаних освітніх середовищ та зорієнтовані на розширення пізнавальних, практичних і рефлексивних можливостей навчання. Такий підхід дозволяє акцентувати не стільки на технологічній природі імерсивних рішень, скільки на їх педагогічній функції --- організації навчальної діяльності в умовах занурення, активної взаємодії та контекстуалізованого засвоєння змісту.

Використання міждисциплінарного понятійного апарату, що поєднує педагогічний, психологічний і технологічний дискурси, зумовлює потребу в чіткому розмежуванні понять «імерсивні технології», «імерсивні освітні ресурси», «цифровий дизайн», «професійна готовність» та «педагогічна ефективність», кожне з яких виконує специфічну функцію у структурі досліджуваного процесу.

Наголошено, що імерсивність у педагогічному вимірі не зводиться до ефекту присутності чи посиленої візуалізації навчального матеріалу. Її сутність пов'язана з потенціалом трансформації логіки навчальної діяльності, характеру взаємодії суб'єктів освітнього процесу та способів конструювання навчального досвіду. Імерсивні освітні ресурси створюють умови для переходу від репродуктивних моделей навчання до діяльнісно-орієнтованих і проєктних форм роботи, у межах яких здобувач освіти виступає активним учасником навчального процесу, а навчальний контент набуває функції середовища для дії, дослідження й рефлексії.

Звернення до наукових джерел переконує: педагогічний потенціал імерсивних технологій реалізується не автоматично, а вимагає системного підходу, який поєднує технологічні можливості з педагогічним проєктуванням, методичною підготовкою вчителя та усвідомленням можливих ризиків. Саме така збалансована позиція забезпечує перехід від декларативного впровадження імерсивних технологій до їх педагогічно ефективного використання.

Систематизація вітчизняного досвіду використання імерсивних технологій у підготовці майбутніх учителів інформатики засвідчила фрагментарний характер їх інтеграції в освітньо-професійні програми. Порівняльний аналіз освітньо-професійних програм показав, що вітчизняна педагогічна освіта перебуває переважно на етапі накопичення цифрово-технологічних передумов для впровадження імерсивних освітніх ресурсів, не забезпечуючи їх цілісного педагогічного осмислення та концептуального закріплення в структурі професійної підготовки. У більшості проаналізованих програм цифрові технології декларуються як важливий компонент підготовки, водночас імерсивна складова не визначається як самостійний об'єкт навчання, окрема освітня мета або структурно окреслений елемент освітньої траєкторії майбутнього вчителя інформатики.

Встановлено, що ідеї імерсивності представлені переважно імпліцитно - через формулювання, пов'язані із застосуванням інформаційно-комунікаційних технологій, цифрових освітніх ресурсів, комп'ютерного моделювання, візуалізації та програмування. Такі результати створюють функціональні передумови для потенційного використання імерсивних освітніх ресурсів, однак не забезпечують їх педагогічної, дидактичної та методичної конкретизації як окремої складової професійної компетентності майбутнього вчителя. Унаслідок цього імерсивні технології залишаються поза логікою цілеспрямованого формування професійних дій і педагогічних рішень.

Аналіз сукупності навчальних дисциплін засвідчив, що потенціал інтеграції імерсивних освітніх ресурсів зосереджується насамперед у курсах комп'ютерної графіки, комп'ютерного моделювання, мультимедійних технологій, STEM/STEAM-освіти, робототехніки, вебтехнологій, а також у вибіркових дисциплінах інноваційного спрямування. Водночас така інтеграція має несистемний характер і не супроводжується чітким визначенням дидактичних цілей, методів навчання, логіки поетапного формування відповідних умінь і критеріїв оцінювання результатів. У межах зазначених дисциплін імерсивні освітні ресурси постають радше як технологічна можливість, ніж як педагогічно осмислений інструмент навчання.

У підсумку імерсивні освітні ресурси використовуються переважно на інструментально-демонстраційному рівні або як допоміжний компонент окремих дисциплін, проєктних завдань і практик, не набуваючи статусу самостійного об'єкта навчання та системоутворювального чинника професійної підготовки майбутніх учителів інформатики.

Результати теоретичного аналізу та узагальнення вітчизняного досвіду засвідчують, що навчання розробки імерсивних освітніх ресурсів у підготовці майбутніх учителів інформатики набуває особливої актуальності в умовах цифрової трансформації освіти. Імерсивні технології демонструють значний педагогічний потенціал, однак його реалізація у вітчизняній педагогічній освіті має переважно фрагментарний характер і не супроводжується цілісною методичною концепцією. Виявлені суперечності між задекларованою значущістю цифрових та імерсивних технологій і відсутністю системного підходу до формування відповідних професійних умінь зумовлюють потребу в спеціально організованому педагогічному впливі.

Методика навчання розробки імерсивних освітніх ресурсів постає не як додатковий або факультативний компонент професійної підготовки, а як необхідна умова формування готовності майбутнього вчителя інформатики до проєктування педагогічно доцільних, безпечних і дидактично обґрунтованих цифрових середовищ. Її значущість визначається здатністю забезпечити перехід від інструментального використання імерсивних технологій до їх усвідомленої інтеграції в логіку навчального процесу, орієнтованого на розвиток пізнавальної активності, практичних умінь і рефлексивного досвіду здобувачів освіти.

Виявлені теоретичні положення та результати аналізу освітніх програм створюють підґрунтя для подальшого обґрунтування методичного забезпечення навчання розробки імерсивних освітніх ресурсів, визначення педагогічних умов, методів, форм і засобів організації цього процесу у закладах вищої педагогічної освіти.

\textbf{\hfill\break
}

\textbf{РОЗДІЛ 2.}

\textbf{МЕТОДИЧНЕ ЗАБЕЗПЕЧЕННЯ НАВЧАННЯ РОЗРОБКИ}

\textbf{ІМЕРСИВНИХ ОСВІТНІХ РЕСУРСІВ}

У розділі розкрито організаційно-педагогічні засади навчання майбутніх учителів інформатики розробки імерсивних освітніх ресурсів у контексті цифрової трансформації освіти. Проаналізовано наукові підходи та принципи, які визначають методологічні засади підготовки здобувачів до створення VR/AR-рішень і моделювання інтерактивних навчальних середовищ. Запропоновано модель формування вмінь створювати імерсивні освітні ресурси у процесі фахової підготовки. Окреслено педагогічні умови, що забезпечують ефективність цього процесу.

\textbf{2.1.~Дидактичні підходи та принципи навчання майбутніх учителів інформатики розробки імерсивних освітніх ресурсів.}

Розвиток імерсивних технологій відкриває нові можливості для цифрової трансформації педагогічної освіти, зокрема у сфері підготовки майбутніх учителів інформатики. Сучасна педагогічна парадигма орієнтує освітній процес на формування не лише цифрових, а й проєктно-технологічних компетентностей, що дозволяють створювати інноваційні освітні продукти у форматі VR/AR/MR. Осмислення методологічних основ навчання студентів розробки імерсивних освітніх ресурсів потребує визначення комплексу наукових підходів, здатних забезпечити наукову цілісність і практичну ефективність підготовки.

У своєму розумінні дидактичних підходів та принципів навчання майбутніх учителів виходимо з сучасного розуміння понять «дидактичний підхід» та «принципи навчання». У науково-педагогічному дискурсі дослідники пропонують широкий перелік дидактичних підходів, які відображають різні методологічні та методичні орієнтири професійної підготовки. Назвемо деякі з них, які, на наш погляд, більш дотичні до нашої теми. Так, \textbf{В.~Радкевич} звертає увагу на значущість знаннєво орієнтованого, особистісно зорієнтованого, діяльнісно-розвивального та компетентнісного підходів~\hl{{[}Р-2{]}.} \textbf{Л.~Ємчик} виокремлює системний, прогностичний, екологічний та компетентнісний підходи, підкреслюючи їхню здатність забезпечувати цілісність і випереджувальність освітнього процесу~\hl{{[}Є-1{]}}. \textbf{А.~Коробко} наголошує на важливості ціннісно орієнтованого, проблемно орієнтованого, аксіологічного та культурологічного підходів як підґрунтя розвитку професійної ідентичності та світоглядної культури майбутнього педагога~\hl{{[}К-16{]}.} \textbf{Л.~Цвіркун,} у свою чергу, акцентує на компетентнісному, системному та особистісно зорієнтованому підходах, розглядаючи їх як ключові в забезпеченні ефективності підготовки фахівців\hl{~{[}Ц-2{]}. \textbf{Н.~Гончарова} зосереджується на обґрунтуванні можливостей використання диференційованого підходу у доборі методичних матеріалів~{[}Г-17{]}.}

З огляду на значну кількість дидактичних підходів, репрезентованих у сучасній науково-педагогічній літературі, а також з урахуванням специфіки предмета й об'єкта нашого дослідження було виокремлено низку концептуально важливих підходів -- аксіологічний, діяльнісний, інноваційний, інтегративний, когнітивний, компетентнісний, конструктивістський, контекстний, рефлексивний, синергетичний і системний, -- які, на нашу думку, є доцільними для врахування під час організації навчання розробки імерсивних освітніх ресурсів майбутніми учителями інформатики.

Розглянемо окреслені дидактичні підходи більш детально, оскільки вони формують теоретичне підґрунтя ефективної підготовки майбутніх учителів інформатики до розробки ІОР.

\emph{\hl{Аксіологічний підхід}} «полягає в гуманістичній спрямованості освітнього процесу на розвиток особистості як мету, результат і головний критерій її ефективності»~\hl{{[}Б-10, c.11{]}}. Звертаємо увагу і на твердження про те, що «аксіологічний підхід визнає особистісно значущі та життєво важливі потреби людини як найвищу цінність»~\hl{{[}К-21, c.49{]}}, орієнтуючи освітній процес на саморозвиток суб'єктів і запобігаючи формуванню інтуїтивних уявлень про цінності.

У підготовці майбутніх учителів інформатики до розробки ІОР аксіологічний підхід підкреслює ціннісний вимір професійної діяльності, зокрема відповідальність за створення цифрового освітнього контенту, що відповідає етичним нормам (коректність подання інформації, уникнення маніпулятивних елементів, безпечна взаємодія користувачів із цифровим середовищем), принципам гуманізму, поваги до особистості, відповідального використання технологій та академічної доброчесності. Важливим стає усвідомлення здобувачами наслідків застосування створених ними VR/AR-ресурсів у реальному освітньому середовищі.

\emph{\textbf{Дяльнісний підхід}} передусім пов'язаний із організацією діяльності, яка забезпечує розвиток і становлення особистості майбутнього вчителя інформатики \hl{{[}Б-10, с. 11{]},} та визначає орієнтири його практичної підготовки до розробки ІОР. Підхід акцентує увагу на поступовому формуванні навичок через виконання практикоорієнтованих завдань, у яких здобувач виступає активним суб'єктом створення імерсивного контенту. Замість пасивного сприйняття інформації, здобувачі активно залучаються до процесу розробки VR/AR-рішень, експериментують з різними інструментами та технологіями, вирішують професійні проблеми та приймають рішення, що сприяє грунтовному засвоєнню матеріалу та розвитку практичних навичок.

\hl{\emph{\textbf{Інноваційний підхід}}} забезпечує відкритість системи підготовки до технологічних змін, інтеграцію новітніх інструментів і платформ у процес навчання. У наукових джерелах інноваційний підхід трактується як «єдиний змістовий комплекс принципів і методів, що використовуються в освітньому процесі та спрямовані на зміну вимог до якості освіти»\hl{~{[}С-18, с.~33{]}. Підхід зумовлює необхідність гнучкої й адаптивної організації підготовки, здатної оперативно інтегрувати новітні інструменти, платформи та технологічні рішення. У підготовці майбутніх учителів інформатики інноваційний підхід орієнтує на формування готовності працювати з новими цифровими засобами, оскільки передбачає «впровадження в освітній процес інформаційно-комп'ютерних технологій, інноваційних форм і методів навчання»~{[}С-18, с.34{]}, що створює умови для практичного опанування сучасних VR/AR-платформ, конструкторів, програмних середовищ та їх використання під час розробки імерсивних освітніх продуктів.}

\emph{Інтегративний підхід} забезпечує цілісність підготовки майбутніх учителів інформатики до розробки ІОР, дозволяє поєднати технічну, дидактичну та психологічну компоненти підготовки в єдину систему. Створення ІОР вимагає не лише технічних знань, але й розуміння педагогічних принципів та психологічних особливостей учнів. Інтегративний підхід забезпечує гармонійне поєднання цих компонентів, коли майбутній учитель одночасно опановує технологічні інструменти, методичні принципи та особливості сприймання інформації учнями, що дозволяє створювати ефективні та цікаві навчальні матеріали.

На думку сучасних вчених \emph{«когнітивний підхід} передбачає орієнтацію не на пасивне засвоєння готових знань чи відтворення способів вирішення завдань, а на активну розумову діяльність, аналіз, оцінювання, прогнозування і рефлексію здобувачів освіти у процесі професійної підготовки»~\hl{{[}Р-3, с.86{]}.} Когнітивний підхід в аспекті нашої теми у робить акцент на процесах сприймання, осмислення й опрацювання інформації, що реалізується через використання візуалізації, моделювання та симуляцій як засобів поглиблення розуміння навчального матеріалу. Використання візуальних образів, моделей та симуляцій дозволяє зробити навчання більш наочним та зрозумілим, що сприяє кращому засвоєнню матеріалу.

\emph{\hl{Компетентнісний підхід}} забезпечує спрямованість підготовки на формування інтегрованих професійних умінь~-- проєктування, програмування, педагогічного дизайну та цифрової креативності. За таких умов здобувачі мають не лише опановувати теоретичні знання, а й набувати здатності до їх практичного застосування у процесі створення реальних імерсивних освітніх продуктів, що відповідає розумінню компетентності як \textbf{«}спроможності кваліфіковано здійснювати діяльність, виконувати завдання або роботу, поєднуючи знання, уміння, навички та ставлення»~\hl{{[}Д-7, с.240{]}.} У межах компетентнісного підходу особливого значення набуває оволодіння навичками проєктування сценаріїв, інтерактивного програмування, розробки візуального дизайну та креативного вирішення педагогічних завдань, оскільки професійна компетентність трактується як «інтегративна якість особистості, що виявляється у здатності ефективно діяти на основі практичного досвіду, знань і вмінь»~\hl{{[}Д-7, с.96{]}.}

\emph{Конструктивістський підхід} передбачає активне залучення здобувачів до побудови власного знання через створення цифрових освітніх продуктів і взаємодію у віртуальному середовищі. Ми зважаємо на думку В.~Соловйова із співавторами, про те, що «конструктивізм виходить з того, що навчання~-- це активний процес, у ході якого суб'єкти активно конструюють знання на основі власного досвіду. Ідеї конструктивізму виражені й у теорії діяльності, згідно з якою діяльність і дії є основою психічного розвитку»~{[}\hl{С-19, с.165}{]}. Тобто здобувачі не просто отримують готові знання, а й самостійно визначають логіку, структуру та функціональне призначення імерсивних ресурсів, активно їх конструюють їх.

\emph{Контекстний підхід} спрямовує навчання у професійне поле, забезпечуючи його максимальну наближеність до реальних умов педагогічної діяльності та відтворення ситуацій шкільної практики, у межах яких ІОР набуває конкретного функціонального призначення. Така орієнтація відповідає сутності контекстного навчання, яке, за визначенням С.~Скворцової, «проектує освітній процес у вищому навчальному закладі як максимально наближений до майбутньої професійної діяльності»~\hl{{[}С-20, с.66{]}. Це також сприяє} розумінню професійної підготовки вчителя як процесу формування здатності до усвідомленого виконання педагогічних дій у реальних освітніх умовах, адже, за зауваженням Т.~Скорик, «підготовка майбутнього вчителя має забезпечувати «якісне виконання педагогічних дій» через формування загальних і спеціальних компетентностей~\hl{{[}С-21, с.64{]}.}

У процесі підготовки майбутні вчителі інформатики залучаються до моделювання педагогічних ситуацій, що дає змогу усвідомити, як створені VR/AR-рішення використовуватимуться в реальному навчальному середовищі, з якими освітніми завданнями вони співвідносяться та які педагогічні ефекти мають забезпечувати. Саме така організація навчання забезпечує трансформацію навчальної діяльності у професійну, оскільки «контекстним навчанням є навчання, в якому динамічно моделюється предметний та соціальний зміст професійної діяльності»\hl{~{[}С-20, с.67{]}.} Така організація навчання також актуалізує практичну значущість імерсивних ресурсів, орієнтує здобувачів на аналіз умов їх застосування, специфіки шкільного освітнього простору й вікових особливостей учнів, а також стимулює адаптацію технологічних і дидактичних рішень до потреб конкретної аудиторії. Загалом контекстний підхід у підготовці майбутніх учителів інформатики до розробки ІОР сприяє формуванню професійної готовності до педагогічно виваженого використання VR/AR-технологій у шкільному навчанні, оскільки забезпечує інтеграцію особистісного розвитку з професійною діяльністю, що, за узагальненням Т.~Скорик, є сутнісною характеристикою професійної стабільності й професіоналізму вчителя\hl{~{[}С-21,~с.65{]}}

\emph{Рефлексивний підхід} спрямований на формування здатності усвідомлювати власну діяльність у процесі розробки ІОР, критично оцінювати її результати та визначати напрями подальшого професійного вдосконалення. У науково-педагогічному дискурсі рефлексія розглядається як психічна діяльність, спрямована на осмислення власних дій, станів і способів діяльності, що забезпечує самопізнання та свідоме управління навчально-професійним розвитком~\hl{{[}У-3, с.506{]}.} Застосування рефлексивного підходу сприяє усвідомленню здобувачем логіки та результативності власних дій на всіх етапах створення імерсивних продуктів, формує здатність аргументовано оцінювати обрані технічні й дидактичні рішення та співвідносити їх із поставленими освітніми цілями. Систематичний аналіз сильних і слабких сторін розроблених VR/AR-рішень, типових проєктних помилок і обмежень стимулює розвиток професійно значущих умінь самокорекції та відповідального ставлення до якості освітнього продукту. Такий підхід узгоджується з позицією В.~Перегудової, яка розглядає педагогічну рефлексію як «систематичний та глибокий процес аналізу й самооцінки педагогічної діяльності», спрямований на вивчення результатів власної практики та пошук шляхів її вдосконалення~\hl{{[}П-13, с.115{]}.} Водночас рефлексивний підхід створює підґрунтя для планування подальших кроків з удосконалення імерсивних продуктів, оптимізації індивідуального проєктного процесу та вибудовування навчально-професійної траєкторії майбутнього вчителя інформатики. За В.~Перегудовою, рефлексія дозволяє майбутнім фахівцям «систематизувати власний досвід та аналізувати його в контексті непередбачуваних обставин», стимулюючи пошук ефективних рішень і впровадження інноваційних підходів \hl{{[}П-13, с. 116{]}.} Відтак \hl{рефлексивний підхід у підготовці майбутніх учителів інформатики до розробки ІОР виступає методологічним підґрунтям формування здатності до самостійного аналізу, самоконтролю та прогнозування власного професійного розвитку, забезпечуючи зростання професійної готовності до педагогічно виваженого використання VR/AR-технологій у сучасному шкільному навчанні.}

У наукових джерелах підкреслюється: синергетика міждисциплінарний напрям, який «розглядає освітні системи як відкриті, здатні до самокоригування та пошуку альтернативних шляхів розвитку»~{[}\hl{К-21, с.49{]}}, що зумовлює переосмислення пізнавального процесу в педагогічній освіті. \emph{Синергетичний підхід} розкриває процес навчання як самоорганізовану систему («з позиції відкритості, співтворчості та орієнтації на саморозвиток, сприяє розв'язанню протиріч, що виникають на різних рівнях освіти»~{[}\hl{Б-10, с.11{]}}), у якій нові знання виникають завдяки взаємодії здобувача, технології та освітнього середовища.

У межах імерсивного навчання синергетичні ідеї набувають особливої значущості, адже нові знання формуються не ізольовано, а в результаті взаємодії здобувача освіти, цифрових технологій та освітнього середовища, у якому технологія постає активним учасником освітньої взаємодії. Така взаємодія сприяє виникненню індивідуальних освітніх траєкторій і стимулює саморозвиток здобувача в умовах відкритої цифрової системи. Відповідно синергетичний підхід у підготовці майбутніх учителів інформатики до розробки ІОР дозволяє розглядати навчання розробки VR/AR-рішень як відкриту самоорганізовану систему, у якій поєднання технологічних, когнітивних і дидактичних компонентів визначає логіку проєктування, вибір методів роботи та практичних інструментів.

\hl{Науковці практично одностайні в тому, що} системний підхід \hl{передбачає розгляд об'єкта як «цілісної системи з визначенням її елементів та упорядкуванням зв'язків між ними»~{[}К-21, с.46{]}, водночас акцентуючи на необхідності виокремлення системотвірних взаємодій і врахування впливу надсистеми та внутрішньої взаємодії підсистем. Така позиція, обґрунтована у працях А.~Коломієць і Н.~Лазаренко, дає підстави розглядати системний підхід як інструмент виявлення причинно-наслідкових залежностей і наукового узагальнення.}

У межах дослідження системний підхід дає змогу розглядати процес розробки ІОР як цілісну педагогічну систему зі стійкими структурними й функціональними зв'язками, у якій усі компоненти перебувають у взаємозумовленій взаємодії. Його застосування уможливлює побудову моделі навчання, у межах якої мета, зміст, методи, засоби та результат поєднані в логічно впорядковану структуру, що забезпечує узгодженість освітніх впливів, оптимізацію навчального процесу та підвищення його ефективності й результативності. У такій логіці системний підхід виступає одним із провідних методологічних засобів організації підготовки майбутніх учителів інформатики до розробки ІОР, оскільки сприяє комплексному опануванню технологій створення VR/AR-матеріалів як цілісної системи, що відображає взаємопов'язаність програмних, дидактичних і дизайнерських компонентів професійної діяльності. Системний підхід також орієнтує на свідомий добір оптимальних способів розв'язання навчально-пізнавальних і проєктних завдань, визначення адекватних теоретичних і практичних засобів підтримки розробки ІОР, у тому числі з використанням евристичних методів пошуку та творчого конструювання, що посилює методичну вивіреність і педагогічну доцільність створюваних VR/AR-рішень.

Функціональне призначення використання тих чи ти підходів у розробці методики навчання майбутніх учителів інформатики розробки ІОР унаочнено на Рис.~2.1.

\includegraphics[width=6.49792in,height=4.46317in]{thesis_final/images/media/image9.png}

\emph{Рис.~2.1.} Дидактичні підходи до формування методики навчання розробки імерсивних освітніх ресурсів

Сукупність поданих \emph{дидактичних підходів} утворює методологічну основу підготовки майбутніх учителів інформатики до розробки імерсивних освітніх ресурсів. Застосування сукупності підходів формує новий тип педагога~-- інженера освітніх технологій, здатного інтегрувати технологічні інновації з педагогічною доцільністю й гуманістичними цінностями освіти.

Ми поділяємо позицію вчених, які стверджують, що реалізація будь-якого дидактичного підходу передбачає добір відповідних принципів навчання~\hl{{[}Ц-2, с.299{]}.} У цьому контексті зосередимо увагу на принципах, які корелюють із визначеними вище дидактичними підходами та формують методологічне підґрунтя навчання майбутніх учителів інформатики розробки імерсивних освітніх ресурсів.

Аналізуючи структуру та функції дидактичних принципів у професійній підготовці фахівців, зокрема педагогів, науковці акцентують на їхньому значенні для забезпечення цілісності, системності й ефективності освітнього процесу. Представлені в сучасній літературі підходи до трактування дидактичних принципів дозволяють окреслити методологічне поле нашого дослідження, а також визначити ті з них, що найбільш повно відповідають специфіці підготовки здобувачів до розробки імерсивних освітніх ресурсів.

У праці Л.~Цвіркун~\hl{{[}Ц-2{]}} дидактичні принципи інтерпретуються як методологічні \emph{засади}, що визначають зміст і логіку професійної підготовки. Дослідниця виокремлює \emph{загальнодидактичні принципи} (науковість, наочність, послідовність та систематичність, зв'язок теорії з практикою, самостійності та активності), а також с\emph{пецифічні для фахової підготовки принципи} (проблемності та професійної мобільності), реалізація яких спрямована на формування професійних умінь через ускладнення навчальних завдань, активізацію мислення та орієнтацію на реальні професійні ситуації. У близькому методологічному ключі О.~Сергійчук розглядає принципи професійної підготовки майбутніх педагогів засобами ІКТ, акцентуючи увагу на інтерактивності, адаптивності, гнучкості, мультимедійності та співпраці \hl{{[}С-22, с. 321--322{]}.}

Позицію \emph{системності} й \emph{наукової цілісності} у трактуванні дидактичних принципів розвивають М.~Гребенюк і В.~Крикун \textbf{{[}Г-16{]},} які розглядають принципи як структурні елементи моделі професійної підготовки, здатні забезпечити узгодженість між цілями, змістом, методами й формами навчання. У межах системного підходу автори виокремлюють принципи \emph{науковості, системності, послідовності, безперервності, автономності навчання та зв'язку з практикою.} Їхній аналіз акцентує на необхідності глибокого розуміння логіки професійної діяльності, адаптивності освітнього процесу до реальних потреб майбутньої роботи та інтеграції сучасних методичних і технологічних рішень. Широкий спектр дидактичних принципів навчання (\emph{систематичність та послідовність, науковість, професійне спрямування, наочність, посильність, міцність, свідомість, активність, виховне навчання, доступність, врахування індивідуальних особливостей та міжкультурної взаємодії)} у їхньому зв\textquotesingle язку із міжпредметним підходом розкриває Ю.~Олізько~\hl{{[}О-5{]}.}

Звернення до наукового доробку вчених дозволяє окреслити базовий \hl{перелік} дидактичних принципів, які забезпечують методологічну цілісність професійної підготовки. Водночас специфіка навчання майбутніх учителів інформатики розробки ІОР потребує уточнення цього переліку з урахуванням технологічної, когнітивної та проєктної складових діяльності. Тому далі розглянемо групу принципів, що, на нашу думку, найбільш повно відповідають завданням підготовки здобувачів до створення VR/AR-ресурсів та забезпечують результативність цього процесу.

Визначення принципів навчання спирається на взаємозв'язок методологічних підходів, що покладений в основу формування цілісної системи методичних засад підготовки майбутніх учителів інформатики до розробки імерсивних освітніх ресурсів. Специфіка цього процесу зумовила виокремлення групи принципів \emph{інтерактивності, науковості, практичної спрямованості, креативності, варіативності, інтеграції, рефлексивності} та \emph{безперервності}, сукупність яких забезпечують якість та результативність процесу навчання.

\emph{Принцип інтерактивності} орієнтує на активну діяльність студента як творця цифрового освітнього контенту. Взаємодія з інструментами VR/AR-розробки відбувається з іншими учасниками навчального процесу. Вихідним положенням якого стає посилення мотиваційного потенціалу~\hl{{[}С-22{]}.} О.~Коваленко та співавтори зазначають, що саме інтерактивність імерсивних середовищ сприяє глибшому залученню студентів до професійної діяльності~\hl{{[}К-22, с.6{]}. С.}~Литвинова разом із співавторами описують створення AR-об'єктів із мультимедійними шарами, що «оживляє» підручник через ігри, квести та симуляції~\hl{{[}Л-12, с.56--57{]}.}

\emph{Принцип науковості} забезпечує обґрунтованість розроблюваних освітніх ресурсів через опору на сучасні положення когнітивної психології, педагогічного дизайну та цифрової педагогіки. Наукова достовірність контенту виступає ключовою умовою якості імерсивних освітніх середовищ, зокрема у процесі моделювання складних явищ і процесів. С.~Литвинова та співавтори зазначають, що VR/AR-технології за дотримання змістових вимог забезпечують відтворення процесів із високим рівнем наукової точності, наголошуючи, що «контент має відповідати цілому набору вимог, у тому числі науковій достовірності»~\hl{{[}Л-12, с.23--24; с.26{]}}. У методичному вимірі М.~Шут і Л.~Благодаренко розглядають принцип науковості як узгодженість змісту навчання з актуальним станом розвитку науки та логікою наукового пізнання, що підвищує концептуальну точність і дидактичну цінність навчального контенту~\hl{{[}Ш-5{]}}.

Принцип \emph{практичної спрямованості} орієнтує підготовку студентів на створення освітніх продуктів, придатних до безпосереднього використання в реальному освітньому середовищі, що передбачає відпрацювання повного циклу їх розробки. О.~Федоренко та О.~Чала акцентують увагу на автоматизації процесів, гейміфікації та прикладному характері імерсивних ресурсів як чинниках посилення практичної орієнтації навчання~\hl{{[}Ф-2, с.7--9{]}.} С.~Литвинова та співавтори зазначають, що використання AR/VR-технологій уможливлює безпечне моделювання потенційно небезпечних хімічних і фізичних процесів~\hl{{[}Л-12, с.18; 55{]},} підкреслюючи, що «за допомогою VR та AR технологій можна провести дослід з небезпечними хімічними речовинами і при цьому не завдати шкоди ні собі, ні оточенню»~\hl{{[}Л-12, с.30{]}.} У цьому контексті О.~Коваленко та співавтори наголошують, що практична спрямованість навчання сприяє формуванню готовності здобувачів освіти до професійної діяльності\hl{~{[}К-22, с.8{]}.}

\emph{Принцип креативності} стимулює пошук нових форм моделювання навчальних ситуацій. Використання інноваційних засобів подання матеріалу зумовило якісне зрушення. Студенти генерують оригінальні ідеї в проектах. С.~Литвинова та співавтори описують процес створення власних AR-об'єктів як потужний інструмент розвитку креативності~\hl{{[}Л-12, с.56--57{]}.} О.~Федоренко та О.~Чала зазначають, що креативність посилюється через взаємодію з інтелектуальними системами та імерсивними середовищами~\hl{{[}Ф-2, с.10--11{]}. Б.~Мінькович обґрунтовує провідну роль інформаційно-комунікаційних технологій у розвитку креативного мислення майбутніх учителів інформатики. Вихідним положенням якого стає застосування евристичних методів (мозковий штурм, синектика та проекти), які уможливлюють колективне генерування ідей~{[}М-8, с.93--95{]}.}

\emph{Принцип варіативності} відкриває можливості використання різних платформ, технологій і форматів створення імерсивного контенту з урахуванням індивідуальних потреб і освітніх запитів здобувачів. С.~Литвинова та співавтори пов'язують реалізацію цього принципу з побудовою індивідуальних освітніх траєкторій і адаптивністю імерсивних ресурсів, наголошуючи, що «індивідуальна освітня траєкторія \ldots{} ґрунтується на виборі здобувачем освіти видів, форм і темпу здобуття освіти»~\hl{{[}Л-12, с.5, 27{]}.} У ширшому педагогічному вимірі О.~Дубасенюк розглядає варіативність як ключову умову реалізації особистісно орієнтованого підходу, що передбачає створення освітнього середовища, здатного забезпечити свободу вибору змісту, форм і способів навчальної діяльності з метою самореалізації та професійного зростання майбутнього вчителя~\hl{{[}Д-8{]}.} У цьому контексті варіативність імерсивних освітніх ресурсів постає не лише технологічною, а й дидактично значущою характеристикою підготовки майбутніх педагогів.

Принцип інтеграції підкреслює необхідність узгодження ІТ-компонентів із педагогічним змістом і психологічними закономірностями навчання, що забезпечує цілісність освітнього процесу. Поєднання VR/AR із традиційними інформаційно-комунікаційними технологіями сприяє глибокій інтеграції цифрових і дидактичних елементів навчання. С.~Литвинова та співавтори визначають «інтеграцію віртуального контенту з фізичним середовищем» як одну з ключових переваг імерсивних технологій~\hl{{[}Л-12, с.6{]}}. У методичному вимірі П.~Мулеса\textbf{~}\hl{{[}М-9{]}} наголошує на принципі інтеграції педагогічних і цифрових технологій як умові створення електронного освітнього середовища, у межах якого поєднуються зміст навчання, інструментарій віртуальної наочності та методи професійної підготовки майбутніх учителів. Концептуально споріднену позицію обґрунтовує Р.~Гуревич\textbf{~}\hl{{[}Г-20{]}}\textbf{,} розглядаючи інтеграцію як багаторівневий дидактичний принцип, що реалізується на методологічному, дидактичному та прикладному рівнях і забезпечує взаємопроникнення знань, методів і технологій у підготовці вчителя. У методиці навчання розробки імерсивних освітніх ресурсів принцип інтеграції реалізується через поєднання педагогічного дизайну, цифрових інструментів і психолого-дидактичних вимог на всіх етапах проєктування ІОР, що забезпечує узгодженість технологічних рішень із навчальними цілями, змістом і очікуваними результатами освітньої діяльності.

Принцип \emph{рефлексивності} спрямований на розвиток здатності до усвідомленого аналізу власної діяльності та результатів проєктування імерсивних освітніх ресурсів, що забезпечує цілеспрямоване вдосконалення VR/AR-рішень. Ключову роль у цьому процесі відіграє зворотний зв'язок як основа формувального оцінювання та корекції технологічних і дидактичних рішень. С.~Литвинова та співавтори розглядають аналіз характеристик об'єктів, рефлексію й формувальне оцінювання як невід'ємні складники роботи з імерсивними ресурсами, наголошуючи на необхідності «продумати форму зворотного зв'язку з учнями»~\hl{{[}Л-12, с.28; 37{]}.}

У психолого-педагогічному вимірі Г.~Кіт\textbf{~}\hl{{[}К-22{]}} трактує рефлексивність як принцип, що забезпечує здатність майбутнього педагога до самоаналізу, оцінювання досягнень і виявлення проблемних зон діяльності, стимулюючи професійне самовдосконалення. У методичному контексті П.~Мулеса~\hl{{[}М-9{]}} підкреслює необхідність систематичного осмислення результатів використання віртуальної наочності та ефективності цифрових інструментів як умови формування методичних умінь і професійної готовності вчителя.

У межах методики навчання розробки імерсивних освітніх ресурсів принцип рефлексивності реалізується через систематичний самоаналіз проєктних рішень і оцінювання відповідності створених ІОР навчальним цілям та очікуваним результатам.

\emph{Принцип безперервності} відображає потребу в постійному оновленні цифрових і методичних компетентностей майбутнього вчителя в умовах стрімкого розвитку імерсивних технологій. У методиці навчання розробки імерсивних освітніх ресурсів цей принцип реалізується через поєднання формалізованого навчання з самонавчанням, мікрокурсами та поетапною проєктною діяльністю, що забезпечує поступове ускладнення завдань і накопичення професійного досвіду. С.~Литвинова та співавтори підкреслюють, що доступність імерсивних ресурсів у будь-якому місці й часі актуалізує парадигму «навчання впродовж життя»~\hl{{[}Л-12, с.4, 21{]}.} У свою чергу О.~Федоренко та О.~Чала наголошують, що безперервність професійного розвитку є ключовою умовою адаптації педагога до швидких змін цифрових технологій та освітніх середовищ~\hl{{[}Ф-2, с.12{]}.}

Узагальнено сукупність дидактичних підходів та принципів навчання майбутніх учителів інформатики розробки імерсивних освітніх ресурсів представлено на \hl{Рис.~2.2.}

\includegraphics[width=6.49792in,height=3.6148in]{thesis_final/images/media/image10.png}

\emph{Рис.~2.2.} Методологічні засади дослідження

Запропоновані методологічні засади охоплюють одинадцять дидактичних підходів, що відображають багатовимірний характер процесу навчання майбутніх учителів інформатики розробки імерсивних освітніх ресурсів. Їх сукупність фіксує не лише технологічний, а й ціннісно-світоглядний, діяльнісний і системний виміри професійної підготовки. З метою забезпечення цілісності та внутрішньої логіки методологічної основи означені підходи згруповано у чотири кластери, кожен з яких репрезентує домінуючий аспект навчального процесу: ціннісно-світоглядний, процесуально-діяльнісний, технологічно-інноваційний та системно-структурний.

\emph{Перший кластер} -- \emph{ціннісно-світоглядні підходи} -- \emph{об'єднує аксіологічний і компетентнісний підходи}, оскільки обидва орієнтують навчальний процес на формування професійних цінностей, світоглядних орієнтирів та внутрішніх мотивів педагога. Аксіологічний підхід задає нормативно-ціннісні рамки використання цифрових технологій, тоді як компетентнісний визначає безперервну природу професійної підготовки, де знання, уміння, цінності та відповідальність становлять єдину систему, що постійно збагачується. Саме з цього кластера логічно випливають принципи науковості та безперервності, які забезпечують фундаментальність і педагогічну вмотивованість розробки VR/AR-ресурсів.

\emph{Другий кластер} -- \emph{процесуально-діяльнісні підходи} -- \emph{включає діяльнісний, конструктивістський, контекстний і рефлексивний підходи.} Спільним для них є те, що вони визначають логіку організації навчальної діяльності студента, спрямованої на активну взаємодію з цифровим середовищем, осмислене конструювання знань та застосування набутого досвіду в різних контекстах. Діяльнісний підхід акцентує виконання практичних завдань, конструктивістський -- створення власного знання, контекстний -- наближення до реальних педагогічних умов, а рефлексивний -- самоаналіз і самокорекцію. Кластер зумовлює принципи інтерактивності та рефлексивності, що відображають активну позицію здобувача у створенні імерсивного контенту.

\emph{Третій кластер} -- \emph{технологічно-інноваційні підходи} -- \emph{охоплює інноваційний, когнітивний і синергетичний підходи.} Їх об'єднує акцент на використанні сучасних цифрових технологій, оптимізації когнітивної діяльності, впровадженні нових форм і засобів навчання, а також на самоорганізації відкритих цифрових систем. Інноваційний підхід визначає потребу у гнучкому включенні новітніх VR/AR-інструментів, когнітивний -- у врахуванні закономірностей сприймання та мислення, синергетичний -- у взаємодії компонентів навчального середовища, що породжують нові освітні ефекти. Логічними принципами цього кластера є варіативність і креативність, які забезпечують здатність майбутнього вчителя працювати з різними технологічними платформами та створювати оригінальні імерсивні рішення.

\emph{Четвертий кластер} -- \emph{системно-структурні підходи} -- \emph{об'єднує системний та інтегративний підходи,} що визначають архітектуру освітнього процесу. Системний підхід забезпечує узгодженість усіх його елементів, а інтегративний -- поєднання технологічних, дидактичних і психологічних компонентів у єдиний зміст підготовки. Цей кластер формує основу для принципів практичної спрямованості та інтегративності, що характеризують підготовку як процес цілісного та всебічного розвитку професійних компетентностей вчителя.

Згруповані підходи утворюють логічно вмотивовану структуру, з якої природно випливають відповідні принципи навчання. Об'єднання підходів у кластери, на нашу думку, підсилює системність, наочність і методологічну завершеність моделі, дозволяючи представити складний процес підготовки майбутніх учителів інформатики до роботи з імерсивними технологіями у зрозумілій і концептуально обґрунтованій формі.

Узагальнення дидактичних підходів та їх кластеризація дали змогу структурувати методологічні засади навчання розробки імерсивних освітніх ресурсів. Виокремлені чотири кластери підходів визначають ціннісні, діяльнісні, технологічні та системно-структурні виміри підготовки майбутніх учителів інформатики, з яких логічно випливають принципи навчання. Сукупно означені методологічні орієнтири формують концептуальне підґрунтя для подальшого конструювання моделі формування вмінь створювати імерсивні освітні ресурси у процесі фахової підготовки, розроблення якої представлено в наступному підрозділі.

\textbf{2.2.~Структурно-функціональна модель підготовки майбутніх вчителів інформатики до розробки імерсивних освітніх ресурсів}

У сучасному науково-педагогічному дискурсі професійне становлення фахівця осмислюється через побудову структурно впорядкованих моделей, що відображають логіку формування професійних умінь і компетентностей. Моделювання в педагогіці розглядається як загальнонауковий метод пізнання, що забезпечує опосередковане вивчення педагогічних об'єктів шляхом створення моделей і перенесення результатів дослідження з моделі на реальний освітній процес (М.~Воровка~\hl{{[}В-4{]}}, О.~Дубасенюк~\hl{{[}Д-9{]}}, О.~Єжова~\hl{{[}Є-3{]},} C.~Косянчук~\hl{{[}К-20{]}, Є.~Лодатко~{[}Л-13{]}} тощо).

Сучасне осмислення моделювання зумовило увагу до класифікації педагогічних моделей, покладеної в основу системних досліджень, що уможливило виокремлення їх навчальних і науково-дослідницьких типів, де прогностичні варіанти відіграють провідну роль у підготовці фахівців до прогнозованих умов діяльності (О.Єжова \hl{{[}Є-3{]}}). Вивчення наукового доробку вчених дозволило нам зазначити тотожність підходів, де вихідним положенням слугує поділ моделей за сферою застосування, забезпечуючи логічний перехід від констатуючих до імітаційних форм, здебільшого орієнтованих на експериментальне вивчення педагогічних систем. Подібна систематизація, на думку науковців, забезпечує глибше розуміння динаміки професійної освіти.

Вітчизняні педагогічні дослідження пропонують різні варіанти таких моделей, які відрізняються набором структурних компонентів і функціональних зв'язків між ними.

У процесі моделювання професійної підготовки більшість науковців виокремлюють структурні блоки/етапи, що забезпечують цілісність і логічну впорядкованість відповідних моделей. Так, А.~Боровик~\hl{{[}Б-10{]}} у дослідженні, присвяченому формуванню дослідницької компетентності майбутнього вчителя історії, обґрунтовує модель, яка включає чотири взаємопов'язані змістові блоки: цільовий, змістовий, організаційно-технологічний та результативний, кожен з яких виконує окрему функцію у структурі фахової підготовки педагога. За позицією С.~Глікман\textbf{~}\hl{{[}Г-20{]}}\textbf{,} структурно-функціональна модель вибудовується як цілісна система взаємопов'язаних компонентів, до складу якої входять: цільовий, теоретико-методологічний, змістово-процесуальний, діагностичний\emph{.} О.~Саркісова~\hl{{[}С-17{]}} запропонувала дидактичну модель, що складається з чотирьох взаємопов\textquotesingle язаних етапів, а саме: організаційного, змістового, процесуального та оцінно-результативного. У процесі розроблення моделі підготовки майбутніх учителів іноземних мов до проєктного навчання І.~Бей~\hl{{[}Б-9{]}} виокремила чотири взаємопов'язані блоки: цільовий, теоретико-методологічний, організаційно-методичний та діагностико-результативний. Схожої логіки дотримується К.~Копняк~\hl{{[}К-19{]}}, яка у структурно-функціональній моделі визначає п'ять взаємопов'язаних блоків: мотиваційно-цільовий, змістовий, організаційно-процесуальний, контрольно-оціночний та результативний. \hl{Аналізуючи методологічний та критеріальний аспекти моделювання педагогічних систем, М.~Воровка пропонує виділяти мотиваційно-діагностичний, навчально-професійний та рефлексивно-оцінний модулі {[}В-4{]}.}

Опрацювання сучасної науково-педагогічної літератури дало підстави стверджувати, що конструювання структурно-функціональної моделі професійної підготовки має ґрунтуватися на внутрішній узгодженості мети, завдань, принципів, підходів, педагогічних умов та критеріїв її реалізації.

У процесі розроблення моделі підготовки майбутніх учителів інформатики до розробки імерсивних освітніх ресурсів ми спиралися на наявні в сучасному науково-педагогічному дискурсі теоретично обґрунтовані моделі професійної підготовки вчителів інформатики за різними напрямами. Зокрема, у працях Р.~Гуревича обґрунтовано модель підготовки майбутніх учителів інформатики до використання інформаційно-комунікаційних технологій, у межах якої акцентовано взаємозв'язок цільового, змістового, технологічного та результативного компонентів професійної підготов\hl{ки~{[}Г-18 -- Г-20{]}.} У зазначених дослідженнях підкреслюється значення системності, поетапності та методичної доцільності організації освітнього процесу, а також орієнтація на формування професійних умінь, релевантних до реальних умов педагогічної діяльності.

Адаптуючи положення окреслених наукових підходів з урахуванням специфіки імерсивних освітніх ресурсів, їх технологічної складності та педагогічної функціональності, нами запропоновано власну модель підготовки майбутніх учителів інформатики, зорієнтовану на поєднання технологічної, дидактичної та проєктної складових фахової підготовки.

Спираючись на окреслені теоретико-методологічні положення (розділ~2.1.), нами розроблено структурно-функціональну модель підготовки майбутніх учителів інформатики до розробки ІОР у процесі фахової підготовки. Така модель розглядається як цілісна система, у межах якої окремі структурні компоненти та функціональні зв'язки між ними забезпечують цілеспрямованість і логіку процесу формування професійних умінь. У ході її конструювання враховано сукупність ключових структурних компонентів, які в сучасних педагогічних дослідженнях розглядаються як методологічно значущі для формування професійних умінь фахівців означеного профілю та забезпечення їх готовності до педагогічно вмотивованого використання ІТ у майбутній професійній діяльності. Усі складники моделі є взаємопов'язаними, взаємозумовленими й займають визначене місце в загальній структурі фахової підготовки. Виходячи з цього, структурно-функціональна модель ІОР освітніх ресурсів передбачає поетапну організацію освітнього процесу, що охоплює \emph{мотиваційно-цільовий,} \emph{теоретико-методологічний, змістовий, організаційно-технологічний та рефлексивно-оцінний блоки.}

Зазначена структурно-функціональна модель підготовки майбутніх учителів інформатики до розробки імерсивних освітніх ресурсів, що відображає логіку взаємодії її основних блоків, подана на Рисунку 2.3.

\includegraphics[width=6.48958in,height=8.48958in]{thesis_final/images/media/image11.png}

\emph{Рис.~2.3.} Структурно-функціональна модель навчання майбутніх учителів інформатики розробки імерсивних освітніх ресурсів

Подальший виклад спрямовано на розкриття змістового наповнення виокремлених структурних блоків моделі з урахуванням їхньої цільової спрямованості, функціонального призначення та методологічних орієнтирів.

\emph{\textbf{Мотиваційно-цільовий блок}} у структурно-функціональній моделі підготовки майбутніх учителів інформатики до розробки ІОР виконує вихідну, смислотвірну та спрямувальну функції. Блок передбачає обґрунтування соціального запиту на підготовку професійно компетентних майбутніх учителів інформатики, здатних до розробки та педагогічно доцільного використання ІОР, врахування потреб особистості у професійному саморозвитку, цифровій мобільності та адаптації до умов трансформації освітнього середовища. Реалізація зазначених положень здійснюється з урахуванням особливостей, принципів і завдань навчання, що зумовлюються метою моделі - формування готовності майбутніх учителів інформатики до розробки імерсивних освітніх ресурсів у межах фахової підготовки.

Нормативно-правове підґрунтя розробленої моделі формується на основі положень Закону України «Про вищу освіту», проєкту стандарту вищої освіти за спеціальністю 014.09 «Середня освіта (Інформатика)» для першого (бакалаврського) рівня, а також чинних ОПП і робочих навчальних планів підготовки майбутніх учителів інформатики, що впроваджуються в закладах вищої освіти у межах відповідних освітніх програм бакалаврського рівня.

Саме в межах цього блоку закладаються ціннісно-мотиваційні основи професійної діяльності, формується усвідомлене ставлення здобувачів освіти до ІТ як до важливого ресурсу сучасної педагогічної практики, а також визначається загальне спрямування подальших етапів підготовки.

\emph{Метою} мотиваційно-цільового блоку є формування готовності майбутніх учителів інформатики до розробки і педагогічно доцільного використання ІОР у професійній діяльності. Мета постає не лише як очікуваний результат навчання, а як інтегративний орієнтир, що поєднує професійні, технологічні та педагогічні аспекти підготовки, узгоджуючи їх із сучасними вимогами до цифрової трансформації освіти.

До основних \emph{завдань} мотиваційно-цільового блоку належить розвиток стійкої мотивації майбутніх учителів інформатики до опанування ІТ як складника сучасної педагогічної практики; забезпечення орієнтації здобувачів освіти на досягнення очікуваних результатів професійної підготовки, пов'язаних із проєктуванням та реалізацією ІОР; формування професійної рефлексії щодо власної діяльності у сфері створення та застосування ІОР. Реалізація цих завдань сприяє усвідомленню значущості ІТ для підвищення якості навчання та розвитку інноваційного педагогічного мислення.

\emph{\textbf{Теоретико-методологічний блок}} у структурно-функціональній моделі підготовки майбутніх учителів інформатики до розробки ІОР виконує концептуально-обґрунтовувальну функцію. Саме в його межах визначаються наукові підходи та принципи, які задають методологічну рамку всього процесу професійної підготовки майбутніх учителів інформатики й забезпечують його внутрішню логічну узгодженість.

\emph{Метою} теоретико-методологічного блоку є формування цілісного науково-педагогічного підґрунтя підготовки майбутніх учителів інформатики, що забезпечує осмислене поєднання теоретичних знань про ІОР з методикою їх педагогічного проєктування та використання у професійній діяльності. У цьому контексті теоретико-методологічні засади виступають джерелом знань та орієнтиром для вибору адекватних способів організації навчання, форм і методів освітньої взаємодії.

До основних \emph{завдань} теоретико-методологічного блоку належить обґрунтування сукупності наукових підходів, що забезпечують ефективність підготовки майбутніх учителів інформатики до розробки ІОР; визначення принципів навчання, які відображають специфіку ІТ та педагогічної діяльності в умовах цифрової трансформації освіти; інтеграція положень педагогіки, психології, інформатики та цифрових технологій у єдину методологічну систему. Реалізація цих завдань сприяє формуванню методологічної культури майбутніх учителів та здатності свідомо застосовувати теоретичні знання у практичній діяльності.

Теоретико-методологічний блок спирається на сукупність дидактичних підходів (докладно у підрозділі 2.1.), зокрема аксіологічний, діяльнісний, інноваційний, інтерактивний, когнітивний, компетентнісний, конструктивістський, контекстний, рефлексивний, синергетичний і системний, які забезпечують цілісне осмислення процесу підготовки майбутніх учителів інформатики до розробки ІОР та орієнтують його на активну, усвідомлену й професійно значущу діяльність. Застосування зазначених підходів дозволяє поєднати ціннісні орієнтири професійної підготовки з практичним досвідом проєктування імерсивних освітніх рішень, а також врахувати когнітивні, діяльнісні й рефлексивні аспекти становлення майбутнього педагога.

Систему принципів навчання у межах теоретико-методологічного блоку становлять принципи науковості, інтерактивності, варіативності, практичної спрямованості, інтеграції, рефлексивності, креативності та безперервності. Їх реалізація забезпечує методологічну обґрунтованість і педагогічну доцільність процесу підготовки, сприяє узгодженню теоретичних положень із практичними завданнями розробки ІОР та створює підґрунтя для послідовного переходу до змістового блоку моделі, зберігаючи наступність між цільовими орієнтирами й практичною підготовкою майбутніх учителів інформатики.

\emph{\textbf{Змістовий блок}} у структурно-функціональній моделі підготовки майбутніх учителів інформатики до розробки ІОР, спрямований на формування цілісного комплексу знань і уявлень, необхідних для усвідомленого проєктування та педагогічно доцільного використання технологій занурення. У цьому блоці конкретизовано навчальний зміст підготовки з урахуванням як технологічних, так і дидактичних вимог до ІОР.

\emph{Метою змістового блоку} є наповнення процесу підготовки структурованими та науково обґрунтованими знаннями про ІОР, принципи їх функціонування, можливості й обмеження використання в освітньому середовищі, а також про засади дидактичного дизайну імерсивних ресурсів. У цьому контексті визначається зміст навчальних дисциплін професійної підготовки, що входять до обов'язкових компонентів ОПП, а також вибіркових дисциплін, орієнтованих на поглиблене вивчення програмування, комп'ютерної графіки, моделювання, інтерактивних технологій, основ VR/AR/MR та дидактичного дизайну цифрових освітніх ресурсів. Така структура змісту забезпечує системне опанування теоретичних основ і технологічних інструментів, необхідних для створення імерсивних освітніх продуктів.

До \emph{основних завдань} змістового блоку належить інтеграція знань з інформатики, педагогіки, психології навчання та освітнього дизайну, формування у майбутніх учителів інформатики цілісного уявлення про імерсивний освітній ресурс як дидактично виважений продукт педагогічної діяльності, а не лише як технічну розробку.

За основу було використано \hl{ОПП Інформатика. програмування КДПУ,} де було виявлені ключові дисципліни циклу психолого-педагогічної, методичної та науково-предметної підготовки, зокрема «Методика навчання інформатики» (ОК 15), «Цифрові технології в освіті» (ОК 18), «Розробка цифрових освітніх ресурсів» (ОК 21), «Комп'ютерна графіка» (ОК 24), «Людино-машинна взаємодія» (ОК 32), а також навчальні й виробничі практики (ОК 36, ОК 37, ОК 40, ОК 42), формують комплекс знань, умінь і навичок, опанування яких забезпечує готовність майбутніх учителів інформатики до розробки імерсивних освітніх ресурсів.

Зазначені освітні компоненти безпосередньо спрямовані на розвиток інформаційно-цифрової компетентності (СК-03), проєктувальної компетентності (СК-09), інноваційної компетентності (СК-13) та здатності проєктувати й розробляти програмне забезпечення з використанням сучасних парадигм програмування (СК-16), що корелює з програмними результатами навчання ПР 02, ПР 03, ПР 09, ПР 13 та ПР 16. У центрі змістового блоку запропонованої структурно-функціональної моделі перебуває дидактичне наповнення й методичне забезпечення спеціально спроєктованого курсу \emph{\textbf{«Розробка імерсивних освітніх ресурсів»}}, орієнтованого на інтеграцію педагогічних, технологічних і дизайнерських компонентів та формування здатності до педагогічно вмотивованого проєктування, створення й використання імерсивних технологій у реальному освітньому середовищі.

Вагоме місце у структурі змістового блоку займає навчально-дослідницька та проєктна діяльність здобувачів освіти, реалізація якої забезпечується в межах фахових і методичних дисциплін, насамперед курсу з методики навчання інформатики, розробки цифрових освітніх ресурсів, комп'ютерної графіки, людино-машинної взаємодії та програмування. У процесі виконання курсових і проєктних завдань студенти застосовують набуті знання й уміння для проєктування та створення власних імерсивних освітніх ресурсів, поєднуючи технологічні рішення з дидактичними цілями навчання. Особливого значення набуває практична підготовка, представлена навчальними та виробничими практиками, зокрема педагогічною та практикою в закладах освіти, під час яких майбутні вчителі інформатики здійснюють апробацію створених імерсивних ресурсів у реальному або наближеному до реального освітньому середовищі, оцінюють їх дидактичну доцільність, функціональність і відповідність освітнім результатам.

Змістовий блок зорієнтований на цілісне поєднання навчального матеріалу з майбутньою професійною діяльністю вчителя інформатики та враховує специфіку розробки імерсивних освітніх ресурсів. Його реалізація спирається на систему дидактичних підходів і принципів, що визначають логіку відбору, структурування й оновлення навчального змісту.

Реалізація змістового блоку створює методичне й дидактичне підґрунтя для переходу до \emph{\textbf{організаційно-технологічного}} блоку моделі, у межах якого засвоєні знання набувають прикладного характеру та трансформуються у практичні вміння і навички розробки, апробації та педагогічного використання імерсивних освітніх ресурсів у реальному або наближеному до реального освітньому середовищі.

\emph{Метою} організаційно-технологічного блоку є забезпечення практико-орієнтованої реалізації процесу підготовки майбутніх учителів інформатики до розробки імерсивних освітніх ресурсів, узгодження цілей і завдань фахової підготовки з вимогами освітніх стандартів, специфікою професійної діяльності вчителя інформатики та технологічними особливостями імерсивних освітніх рішень. У межах цього блоку конкретизуються завдання формування відповідних умінь, визначається місце імерсивних технологій у структурі освітньої програми, а також забезпечується взаємозв'язок між теоретичною підготовкою, практичною діяльністю та проєктною роботою студентів.

Організаційно-технологічний блок охоплює процесуально-діяльнісний складник, спрямовану на безпосереднє формування практичних умінь і досвіду проєктної діяльності майбутніх учителів інформатики. В межах цього блоку відбувається перехід від засвоєння теоретичних положень до їх реалізації у практиці створення імерсивних освітніх продуктів, що відповідають дидактичним цілям, навчальному контексту та потребам здобувачів освіти.

До основних завдань організаційно-технологічного блоку належить організація практико-орієнтованої діяльності студентів, залучення їх до виконання навчально-проєктних завдань, що передбачають застосування інструментів програмування, моделювання, інтерактивної візуалізації, прототипування та тестування імерсивних освітніх ресурсів. Важливого значення набуває розвиток умінь планувати й реалізовувати проєкти, працювати в команді, здійснювати рефлексію результатів власної діяльності, а також створення організаційних і технологічних умов для інтеграції аудиторної роботи з самостійною, командною, проєктною та дослідницькою діяльністю здобувачів освіти.

Організаційно-технологічний блок спрямований на активізацію навчальної діяльності здобувачів освіти та створення ними власного освітнього продукту в процесі фахової підготовки. Його реалізація забезпечує логічний перехід до діагностично-результативного блоку, орієнтованого на оцінювання й узагальнення досягнутих результатів підготовки.

\emph{\textbf{Рефлексивно-оцінний блок}} у структурно-функціональній моделі підготовки майбутніх учителів інформатики до розробки ІОР спрямований на визначення рівня сформованості професійних умінь, узагальнення результатів навчальної діяльності та оцінювання ефективності реалізованої підготовки. На цьому етапі відбувається зіставлення запланованих цілей із досягнутими результатами, що дає змогу виявити ступінь готовності здобувачів до розробки і впровадження ІОР в у педагогічну практику.

Метою оцінно-результативного блоку є комплексне оцінювання результатів формування вмінь створювати імерсивні освітні ресурси з урахуванням технологічних, дидактичних та проєктних складників підготовки. У межах цього блоку визначаються критерії, показники та рівні сформованості відповідних умінь, здійснюється аналіз якості створених освітніх продуктів, а також оцінюється здатність майбутніх учителів інформатики обґрунтовувати педагогічну доцільність використання розроблених ресурсів.

До основних завдань оцінно-результативного етапу належить здійснення поточного й підсумкового контролю, організація самооцінювання та взаємооцінювання, аналіз динаміки професійного зростання студентів. Важливим компонентом цього етапу є рефлексія власної діяльності, що передбачає осмислення отриманого досвіду, виявлення труднощів і визначення шляхів подальшого вдосконалення професійних умінь.

Реалізація цього завдань цього блоку дозволяє не лише зафіксувати досягнуті результати, а й окреслити перспективи подальшого професійного вдосконалення майбутніх учителів інформатики у сфері розробки імерсивних освітніх ресурсів.

Завершальним елементом структурно-функціональної моделі виступає результативний компонент, який відображає ефективність реалізації методики навчання розробки імерсивних освітніх ресурсів у процесі фахової підготовки майбутніх учителів інформатики. Реалізація рефлексивно-оцінного блоку забезпечує систематичне осмислення студентами власного досвіду проєктування та створення імерсивних освітніх ресурсів, співвіднесення проміжних і підсумкових результатів із цілями навчання, а також оцінювання сформованості професійних умінь відповідно до визначених критеріїв і рівнів.

У результаті впровадження запропонованої методики формується готовність майбутніх учителів інформатики до розробки і педагогічно обґрунтованого використання імерсивних освітніх ресурсів, що виявляється у здатності методично доцільно добирати інструменти й технології, проєктувати структуру імерсивного освітнього продукту, інтегрувати його в освітній процес та здійснювати рефлексивне оцінювання власної педагогічної діяльності. Досягнення окресленого результату підтверджує методичну цілісність моделі, узгодженість її структурних блоків і результативність поетапної організації навчання розробки імерсивних освітніх ресурсів.

У межах дослідження означена готовність розглядається як методична характеристика професійної підготовки, що відображає здатність до цілеспрямованого планування, організації та оцінювання процесу створення імерсивних освітніх ресурсів у навчальному середовищі. Виходячи з цього, подальша характеристика готовності здійснюється через виокремлення її структурних компонентів, які забезпечують методичну цілісність і результативність навчання розробки імерсивних освітніх ресурсів у процесі фахової підготовки майбутніх учителів інформатики.

Готовність майбутніх учителів інформатики до розробки й використання імерсивних освітніх ресурсів у межах дослідження трактується як цілісне інтегративне професійно-особистісне утворення, що забезпечує здатність до усвідомленого цілепокладання, педагогічно обґрунтованого проєктування, реалізації та оцінювання результатів освітньої діяльності в цифровому освітньому середовищі. Формування означеної готовності відбувається у процесі поєднання мотиваційних настанов, системи фахових і методичних знань, практичних умінь, технологічних способів дії та рефлексивного осмислення власної педагогічної діяльності.

Визначальну роль у структурі готовності відіграє \emph{\textbf{мотиваційно-ціннісний компонент},} який виконує смислоутворювальну й регулятивну функції та зумовлює спрямованість майбутнього вчителя інформатики на реалізацію методики навчання розробки імерсивних освітніх ресурсів. Його зміст виявляється у сформованості стійких професійних мотивів до педагогічно вмотивованого проєктування імерсивних освітніх продуктів, позитивному ставленні до використання цифрових технологій у навчальному процесі, готовності до безперервного професійного саморозвитку та усвідомленого вибору індивідуальної освітньої траєкторії. Саме цей компонент забезпечує методично обґрунтоване цілепокладання, визначення пріоритетів педагогічної діяльності та внутрішню настанову на досягнення результативності під час розв'язання професійних педагогічних і технологічних завдань.

\emph{\textbf{Когнітивний компонент}} відображає систему професійно значущих знань, необхідних для реалізації методики навчання розробки імерсивних освітніх ресурсів, і формується у процесі опанування змісту фахових, методичних і психолого-педагогічних дисциплін. Він охоплює знання теоретичних засад імерсивного навчання, принципів педагогічного проєктування цифрових освітніх ресурсів, можливостей віртуальних і доповнених середовищ, а також дидактичних вимог до їх використання в освітньому процесі. Когнітивний компонент забезпечує усвідомлене прийняття методичних рішень і перехід від репродуктивного використання готових цифрових продуктів до самостійного педагогічного проєктування.

\emph{\textbf{Діяльнісно-проєктувальний компонент}} характеризує рівень сформованості практичних умінь майбутніх учителів інформатики, пов'язаних із розробкою та впровадженням імерсивних освітніх ресурсів. Його зміст охоплює володіння апаратними й програмними засобами, уміння здійснювати відбір, структурування та педагогічну адаптацію навчального контенту, а також використання спеціалізованих програмних середовищ для створення віртуальних і доповнених освітніх ресурсів. У межах цього компонента формується здатність інтегрувати імерсивні ресурси в структуру уроку або навчального курсу з урахуванням освітніх цілей і методичних вимог.

\emph{\textbf{Технологічний компонент}} розкриває способи організації освітнього процесу з використанням цифрових та імерсивних технологій і відображає готовність до технологічного проєктування навчання. Він передбачає здатність планувати навчальну діяльність, визначати логіку подання навчального матеріалу, організовувати поетапне просування здобувачів освіти до досягнення запланованих результатів, а також добирати адекватні технологічні рішення відповідно до дидактичних завдань. Реалізація цього компонента забезпечує використання імерсивних освітніх ресурсів як інструменту методично виваженої педагогічної діяльності.

\emph{\textbf{Інтерактивний компонент}} відображає здатність майбутнього вчителя інформатики організовувати педагогічну взаємодію в умовах цифрового й імерсивного освітнього середовища. Його зміст пов'язаний із розумінням інтерактивності як методичної характеристики освітнього процесу, що забезпечує активну участь здобувачів освіти, діалог, співпрацю та спільне розв'язання навчальних завдань. У межах використання імерсивних освітніх ресурсів цей компонент реалізується через створення умов для активної пізнавальної діяльності та педагогічно доцільної взаємодії з цифровим середовищем.

\emph{\textbf{Рефлексивно-оцінний компонент}} завершує структуру готовності та забезпечує усвідомлення результатів власної педагогічної діяльності в контексті реалізації методики навчання розробки імерсивних освітніх ресурсів. Він виявляється у здатності оцінювати відповідність створених ресурсів дидактичним цілям, аналізувати ефективність їх використання в освітньому процесі, виявляти проблемні аспекти та здійснювати корекцію педагогічних і технологічних рішень, визначаючи напрями подальшого професійного розвитку.

Компонентна структура готовності майбутніх учителів інформатики до реалізації методики навчання розробки імерсивних освітніх ресурсів узагальнено та подано у \emph{Рисунку 2.4.}

\includegraphics[width=6.49792in,height=2.70034in]{thesis_final/images/media/image12.png}

\emph{Рис.~2.4.} Структура готовності майбутніх учителів інформатики до реалізації методики навчання розробки імерсивних освітніх ресурсів

Представлений рисунок відображає логіку взаємозв'язку мотиваційно-ціннісного, когнітивного, діяльнісно-проєктувального, технологічного, інтерактивного та рефлексивно-оцінного компонентів, які в сукупності забезпечують методичну цілісність процесу підготовки майбутніх учителів інформатики.

Розроблена структурно-функціональна модель підготовки майбутніх учителів інформатики до розробки імерсивних освітніх ресурсів забезпечує цілісне осмислення методики формування відповідних професійних умінь у процесі фахової підготовки. Її реалізація уможливлює системне поєднання цільових орієнтирів, змістових компонентів, організаційно-технологічних механізмів і результативних показників, а також дозволяє окреслити методичні можливості й обмеження організації освітнього процесу. Практичне впровадження моделі створює передумови для цілеспрямованого розвитку здатності майбутніх учителів інформатики до педагогічно вмотивованого та технологічно обґрунтованого проєктування імерсивних освітніх ресурсів у професійній діяльності.

\textbf{2.3. Педагогічні умови та методика навчання студентів розробки імерсивних освітніх ресурсів.}

Подальший хід дослідження зумовив необхідність визначення педагогічних умов, що створюють оптимальні передумови для формування у здобувачів умінь розробляти імерсивні освітні ресурси та інтегрувати їх у навчальне середовище.

Передусім визначимо сутність поняття «педагогічні умови» з огляду на вітчизняні педагогічні розвідки.

\hl{У довідкових джерелах подається низка тлумачень поняття «умова», які, попри варіативність формулювань, мають спільне смислове ядро. Поняття «умова» розуміється як необхідна обставина, що забезпечує можливість здійснення, створення або формування певного явища чи процесу та сприяє досягненню визначеного результату, що зафіксовано у Великому тлумачному словнику сучасної української мови~{[}В-7, с.1506{]}. Розвиваючи загальнофілософське й лексикографічне розуміння зазначеної дефініції, сучасні дослідники конкретизують її у площині педагогічної науки. Зокрема, педагогічні умови трактуються як сприятливі педагогічні обставини та чинники, що свідомо створюються у закладі вищої освіти з метою цілеспрямованого впливу на здобувачів освіти~{[}Є-2, с.32{]}.}

\hl{У науково-педагогічному дискурсі педагогічні умови осмислюються також як узгоджена сукупність внутрішніх параметрів і зовнішніх характеристик організації та функціонування освітнього процесу, здатна забезпечити його високу результативність і відповідність психолого-педагогічним критеріям оптимальності, на чому наголошує В.~Манько~{[}М-10, с.153--161{]}. Концептуально близьку позицію висловлює Б.~Гершунський, який пов'язує педагогічні умови з характеристиками процесу навчання та освітнього середовища, у межах якого формуються й реалізуються педагогічні впливи, спрямовані на досягнення поставлених цілей~{[}Г-21, с.46{]}. Водночас А.~Ашеров і В.~Логвіненко акцентують увагу на тому, що педагогічні умови мають відповідати певним вимогам, серед яких визначальними виступають системний характер, чітка структурованість і наявність стійких зв'язків між структурними компонентами, а також урахування особливостей професійної підготовки здобувачів освіти з позицій їхньої готовності до майбутньої професійної діяльності~{[}А-4{]}.}

\hl{В контексті нашої роботи особливу цікавість викликають праці, де науковці аналізують професійну підготовку майбутніх учителів інформатики та пропонують своє бачення поняття «педагогічні умови» та власний перелік таких умов. Зокрема Т.~Вдовичин зазначає, що «для студентів напряму підготовки «Інформатика» педагогічні умови можна розглядати як певні обставини процесу навчання, в яких, враховано наявні умови навчання, передбачено способи перетворення даних умов у напрямі цілей навчання, відібрано та використано елементи змісту, методи й організаційні форми навчання з урахуванням принципів оптимізації»~{[}В-5, с.227{]}. В контексті формування ІКТ-компетентності С.~Петренко} такі узагальнення: 1)~більшість мовознавців трактують поняття «умова» як обставини реальної дійсності, за яких відбуваються або здійснюються певні події, або як комплекс правил, чи вимог, що встановлені, або встановлюються в певній області життєдіяльності суспільства; 2) філософське тлумачення полягає в тому, що умова -- це те, від чого залежить об'єкт, комплекс об'єктів, характер їх взаємодії, із наявності цієї взаємодії випливає можливість існування, функціонування і розвитку даного об'єкту; 3) педагогічний підхід трактує «умови» як певні обставини, що впливають на появу і розвиток деякого процесу, умови і складають середовище, обставини, в яких вони виникають, існують і розвиваються»~\hl{{[}П-8, с.143{]}.}

Відповідно до напряму досліджень вчені подають різні варіанти «педагогічних умов»:

\begin{itemize}
\item
  «формування мотивації до професійної діяльності шляхом використання VR у процесі професійної підготовки здобувачів вищої освіти» та «впровадження методики використання VR у процесі професійної підготовки здобувачів освіти»~\hl{{[}Є-2, с.32-33{]}};
\item
  «для навчання студентів напряму підготовки «Інформатика»: 1) спрямовання на самореалізацію особистості за допомогою створення та підтримки інновацій в навчальному середовищі; 2) реалізація методичних умінь та навичок у викладачів та забезпечення розвитку їх професійних компетентностей; 3) мотивації до навчальної діяльності та розвиток пізнавальної активності студентів»\hl{~{[}В-5, с.226{]};}
\item
  \hl{«застосування комплексу інноваційних комп'ютерно орієнтованих форм і методів організації навчального процесу студентів; застосування сучасних інформаційно-комунікаційних засобів в різносторонньому і ґрунтовному вивченні математики; реалізація діагностичного підходу до діяльності студентів та її результатів» та «використання універсальних засобів опрацювання інформації з метою розкриття пізнавальних здібностей і творчого потенціалу кожного студента; поєднання дидактичних надбань минулого та сучасних інформаційно-комунікаційних технологій»~{[}П-8, с.143{]}.}
\end{itemize}

Крім того, звертаємо увагу на класифікацію педагогічних умов (психолого-педагогічні, методичні, технологічні, організаційні), запропоновану в наукових джерелах, яку в межах методики навчання розробки імерсивних освітніх ресурсів майбутніх учителів інформатики адаптуємо відповідно до логіки та етапності цього процесу.

Підсумовуючи напрацювання вчених і з урахуванням особливостей предмета нашого дослідження, педагогічні умови підготовки майбутніх учителів інформатики до розробки ІОР доцільно конкретизувати крізь призму наведеної класифікації. У цьому контексті виокремлені нами умови корелюють передусім із психолого-педагогічною та методико-технологічною групами педагогічних умов і водночас відображають їх інтегративний характер.

Психолого-педагогічна умова~-- стимулювання професійної мотивації студентів через організацію навчального процесу, у межах якого створення та педагогічне застосування імерсивних освітніх ресурсів виступає засобом усвідомлення значущості майбутньої професійної діяльності, розвитку рефлексивного ставлення до професії та розкриття дидактичного потенціалу VR/AR-технологій у сучасній школі.

Методична умова -- упровадження цілісної методики навчання, спрямованої на поетапне формування у студентів умінь проєктувати, педагогічно обґрунтовувати й інтегрувати імерсивні освітні ресурси в освітній процес відповідно до цілей і змісту шкільної інформатичної освіти.

Технологічна умова -- забезпечення опанування студентами цифрових інструментів і середовищ для розробки імерсивних освітніх ресурсів, що уможливлює технічну реалізацію VR/AR-проєктів, їх тестування, удосконалення та адаптацію до освітніх потреб здобувачів загальної середньої освіти.

Організаційна умова -- проєктування освітнього процесу на засадах поетапності, міждисциплінарної інтеграції та узгодженості цілей, змісту, методів і засобів навчання, а також визначення критеріїв, показників і рівнів сформованості готовності майбутніх учителів інформатики до розробки імерсивних освітніх ресурсів.

Проведений аналіз обумовив необхідність докладно охарактеризувати означені педагогічні умови, визначивши їхню значущість для ефективної організації підготовки майбутніх вчителів інформатики до створення імерсивних освітніх ресурсів.

\emph{\textbf{Психолого-педагогічна умова,}} що полягає в стимулюванні професійної мотивації студентів через організацію навчального процесу, де створення та педагогічне застосування імерсивних освітніх ресурсів стає потужним засобом усвідомлення значущості майбутньої професійної діяльності, набуває особливої ваги в підготовці майбутніх учителів інформатики. Вихідним положенням такого підходу виступає ідея, що практична діяльність з розробки VR/AR-ресурсів дозволяє студентам відчути реальний вплив цифрових технологій на освітній процес, що зумовило зростання внутрішньої мотивації до професії~\hl{{[}К-23{]}.} Коли майбутні вчителі самі створюють імерсивні середовища для уроків, вони не просто опановують інструменти -- вони переживають, як ці інструменти змінюють сприйняття навчального матеріалу учнями, роблять навчання глибшим і емоційно насиченішим. Саме тому така діяльність уможливила зрушення в професійному самовизначенні студентів, адже вони бачать безпосередній зв'язок між власними зусиллями та потенційними успіхами школярів~\hl{{[}Ре-3, с.112{]}.}

Провідну роль тут відіграє рефлексивне ставлення до професії. Під час створення імерсивних ресурсів студенти змушені постійно аналізувати власні дії: чому саме цей VR-сценарій ефективний, як він впливає на мотивацію учнів, які дидактичні принципи втілено найкраще. Така рефлексія дозволила відзначити істотне посилення усвідомлення професійної ідентичності \hl{}~\hl{{[}О-6{]}.} Студенти починають сприймати себе не як користувачів технологій, а як творців інноваційного освітнього простору, що тотожне переходу від пасивного навчання до активного професійного зростання.

Організація навчального процесу навколо розробки та апробації імерсивних ресурсів забезпечило розкриття дидактичного потенціалу VR/AR-технологій у сучасній школі. Дослідження показують, що імерсивні технології суттєво підвищують залученість учнів, сприяють кращому запам'ятовуванню матеріалу та розвивають критичне мислення\hl{~{[}К-23, с.89--90{]}.} Коли майбутні вчителі інформатики самі створюють такі ресурси, вони не лише вивчають теоретичні аспекти, а й переконуються на практиці в їхній ефективності. Це зумовило глибше розуміння того, як VR/AR може трансформувати традиційні уроки в інтерактивні, мотивуючі досвіди~\hl{{[}П-3{]}.} Такий підхід виявився особливо ефективним у формуванні стійкої професійної мотивації. Студенти, які брали участь у створенні імерсивних проєктів, демонстрували вищий рівень ентузіазму щодо впровадження інновацій у майбутній роботі~\hl{{[}К-24{]}}. Вони усвідомлювали, що володіння цими технологіями робить їх конкурентоспроможними фахівцями, здатними відповідати викликам цифрової освіти.

Отже, організація навчального процесу, де створення та педагогічне застосування імерсивних освітніх ресурсів виступає ключовим елементом, уможливила не лише технічне оволодіння VR/AR, але й глибоке усвідомлення їхньої дидактичної цінності, розвиток рефлексії та стійкої мотивації до професійної діяльності.

\emph{\textbf{Методична умова,}} що полягає в упровадженні цілісної методики навчання, спрямованої на поетапне формування у здобувачів умінь проєктувати, педагогічно обґрунтовувати й інтегрувати імерсивні освітні ресурси в освітній процес відповідно до цілей і змісту шкільної інформатичної освіти, виявилася визначальною для якісної підготовки майбутніх учителів інформатики.

Вихідним положенням розробки методики став принцип системності формування професійних компетентностей. Спочатку здобувачі опановують базові технічні навички створення імерсивного контенту, потім переходять до аналізу дидактичних можливостей VR/AR, а на завершальному етапі вчаться інтегрувати розроблені ресурси в уроки інформатики, чітко узгоджуючи їх із програмними вимогами та віковими особливостями учнів. Така послідовність, на наше переконання, забезпечує уникнення фрагментарного засвоєння знань.

Визначальним елементом цієї методики постає педагогічна аргументація кожного розробленого імерсивного ресурсу. Здобувачі освіти опановують уміння критично осмислювати: яку дидактичну мету реалізує конкретний VR-сценарій, яким чином він забезпечує досягнення запланованих результатів навчання, які способи й прийоми стимулювання пізнавальної активності учнів у ньому закладені.

Цілісна методика уможливлює формування в здобувачів уміння в подальшому інтегрувати імерсивні ресурси саме в шкільну інформатичну освіту; «урізноманітнити форми та засоби роботи зі здобувачами освіти», «підвищити ефективність освітнього процесу інноваційними формами»~\hl{{[}К-25, c.72{]}. За таких умов ІТ імерси}вні технології не просто «прикрашають» урок, а суттєво підвищують його ефективність і відповідність сучасним освітнім стандартам\hl{. О}тже, упровадження цілісної поетапної методики навчання забезпечує не лише оволодіння здобувачами практичними навичками створення імерсивних ресурсів, але й сприяє сформуванню в них здатності свідомо, педагогічно обґрунтовано й системно інтегрувати ці технології в процес шкільної інформатичної освіти.

\emph{Технологічна умова.} Її дотримання спрямоване на формування у майбутніх учителів інформатики здатності комплексно організовувати процес технічного створення імерсивних освітніх продуктів, охоплюючи всі ключові етапи -- від концептуалізації ідеї та цифрового проєктування до програмної реалізації, апробації, аналізу функціонування й подальшого вдосконалення розробленого ресурсу. У межах цієї умови акцент переноситься на роботу з різними імерсивними середовищами, типами інтерфейсів і способами взаємодії користувача з віртуальним контентом. Така підготовка дає змогу студентам усвідомлено добирати й модифікувати технологічні рішення з урахуванням вікових, когнітивних та освітніх особливостей здобувачів загальної середньої освіти.

Водночас технологічна підготовка розглядається не як автономне опанування інструментів, а як складник цілісної професійної підготовки, інтегрований із методичними та педагогічними засадами навчання. Така інтеграція забезпечує готовність майбутнього вчителя не лише використовувати VR/AR-рішення в освітньому процесі, а й гнучко трансформувати їх відповідно до зміни навчальних цілей, змісту теми або умов організації навчання.

\emph{Організаційна умова} передбачає цілеспрямоване проєктування освітнього процесу як логічно впорядкованої та внутрішньо узгодженої системи, у межах якої підготовка майбутніх учителів інформатики до розробки імерсивних освітніх ресурсів здійснюється поетапно, з урахуванням зростання складності навчальних завдань і поступового ускладнення професійних дій. Поетапність організації навчання забезпечує послідовний перехід від засвоєння базових уявлень про імерсивні технології та їх освітній потенціал до інтегрованого застосування знань і вмінь у процесі самостійного проєктування та педагогічного обґрунтування VR/AR-рішень.

Важливим складником цієї умови виступає міждисциплінарна інтеграція, що забезпечує взаємопроникнення змісту фахових, педагогічних і методичних дисциплін. Такий підхід створює підґрунтя для формування цілісного бачення імерсивного освітнього ресурсу як результату поєднання технологічних, дидактичних і психолого-педагогічних компонентів. Унаслідок цього здобувачі освіти усвідомлюють взаємозалежність програмування, цифрового дизайну, методики навчання інформатики та вікових особливостей учнів, що сприяє подоланню фрагментарності підготовки й орієнтації виключно на технічний аспект.

Організаційна умова також передбачає узгодженість цілей, змісту, методів і засобів навчання, що унеможливлює стихійне або епізодичне використання імерсивних технологій у фаховій підготовці. Навчальні цілі конкретизуються через очікувані результати, зміст добирається з урахуванням професійної значущості імерсивних ресурсів для шкільної інформатичної освіти, методи орієнтуються на активну, проєктну та рефлексивну діяльність студентів, а засоби навчання розглядаються як інструменти реалізації цілісного дидактичного задуму. Така узгодженість забезпечує логіку освітнього процесу та сприяє формуванню готовності до системного використання VR/AR-технологій у майбутній професійній діяльності.

Окрему роль у реалізації організаційної умови відіграє визначення чітких критеріїв, показників і рівнів сформованості готовності майбутніх учителів інформатики до розробки імерсивних освітніх ресурсів. Наявність діагностично обґрунтованого інструментарію дозволяє здійснювати моніторинг динаміки професійного становлення студентів, виявляти сильні сторони та проблемні зони підготовки, а також коригувати зміст і методи навчання відповідно до досягнутих результатів. У такий спосіб організаційна умова забезпечує не лише структурну впорядкованість освітнього процесу, а й його результативність, орієнтовану на формування цілісної готовності до педагогічно вмотивованої та методично виваженої розробки імерсивних освітніх ресурсів.

\subsubsection{}\label{section}

\subsubsection{\texorpdfstring{\textbf{Висновки до розділу 2}}{Висновки до розділу 2}}\label{ux432ux438ux441ux43dux43eux432ux43aux438-ux434ux43e-ux440ux43eux437ux434ux456ux43bux443-2}

У розділі обґрунтовано методичне забезпечення навчання майбутніх учителів інформатики розробки імерсивних освітніх ресурсів як цілісну, внутрішньо узгоджену систему, що поєднує методологічні засади, модельну організацію освітнього процесу та сукупність педагогічних умов її ефективної реалізації. Зміст розділу дозволяє розкрити логіку переходу від теоретико-методологічного осмислення проблеми до її практико-орієнтованого розв'язання у межах фахової підготовки.

Встановлено, що навчання розробки імерсивних освітніх ресурсів потребує опори на багатовимірну сукупність дидактичних підходів, які відображають не лише технологічний, а й ціннісно-світоглядний, діяльнісний, когнітивний та системний виміри професійної підготовки. Виокремлені аксіологічний, діяльнісний, інноваційний, інтегративний, когнітивний, компетентнісний, конструктивістський, контекстний, рефлексивний, синергетичний і системний підходи утворюють методологічне підґрунтя формування здатності майбутніх учителів інформатики до педагогічно вмотивованого проєктування імерсивних освітніх продуктів. Їх кластеризація за домінуючими аспектами навчального процесу дозволила структурувати методологічні орієнтири та визначити систему принципів навчання, релевантних специфіці VR/AR-технологій. З'ясовано, що принципи науковості, інтерактивності, практичної спрямованості, інтеграції, креативності, варіативності, рефлексивності та безперервності не лише корелюють із визначеними підходами, а й забезпечують внутрішню логіку методики навчання розробки імерсивних освітніх ресурсів.

Розроблено та теоретично обґрунтовано структурно-функціональну модель підготовки майбутніх учителів інформатики до розробки імерсивних освітніх ресурсів, яка відображає поетапну організацію освітнього процесу та взаємозв'язок його ключових блоків. Визначено, що мотиваційно-цільовий, теоретико-методологічний, змістовий, організаційно-технологічний і рефлексивно-оцінний блоки забезпечують цілісність процесу формування професійних умінь, узгодженість цілей, змісту, методів і результатів підготовки. Запропонована модель дозволяє інтегрувати технологічну, дидактичну та проєктну складові професійної діяльності майбутнього вчителя інформатики, а також забезпечує перехід від засвоєння знань до створення власних імерсивних освітніх продуктів і їх педагогічного обґрунтування. У структурі моделі готовність до розробки імерсивних освітніх ресурсів розглянуто як інтегративне професійно-особистісне утворення, що охоплює мотиваційно-ціннісний, когнітивний, діяльнісно-проєктувальний, технологічний, інтерактивний і рефлексивно-оцінний компоненти.

Визначено й обґрунтовано педагогічні умови, що забезпечують результативність реалізації розробленої моделі та методики навчання розробки імерсивних освітніх ресурсів. Ефективність підготовки майбутніх учителів інформатики зумовлюється цілеспрямованим стимулюванням професійної мотивації студентів через включення їх у діяльність зі створення та педагогічного застосування VR/AR-рішень, упровадженням цілісної поетапної методики навчання, спрямованої на формування здатності проєктувати й методично обґрунтовувати ІОР, забезпеченням технологічної готовності до повного циклу їх розробки, а також організаційним проєктуванням освітнього процесу на засадах поетапності, міждисциплінарної інтеграції та узгодженості цілей, змісту й засобів навчання. Сукупність означених умов створює підґрунтя для системного формування готовності до педагогічно виваженого використання імерсивних технологій у шкільній інформатичній освіті.

Результати другого розділу підтверджують, що методичне забезпечення навчання розробки ІОР має розглядатися як комплексна педагогічна система, у межах якої методологічні підходи, структурно-функціональна модель і педагогічні умови перебувають у взаємозумовленій єдності. Отримані висновки створюють теоретико-методичне підґрунтя для подальшої експериментальної перевірки ефективності розробленої методики та обґрунтування критеріїв і рівнів сформованості готовності майбутніх учителів інформатики до розробки ІОР, що становить зміст наступного розділу дослідження.
